\documentclass[11pt]{article}
\usepackage[T1]{fontenc}
\usepackage[utf8]{inputenc}
\usepackage{lmodern,amsmath,amssymb,amsthm,mathtools}
\usepackage[margin=28mm]{geometry}
\usepackage{microtype,xcolor}
\usepackage[authoryear,round]{natbib}
\usepackage{xurl}
\usepackage[colorlinks=true,allcolors=blue!50!black]{hyperref}
\newtheorem{theorem}{Theorem}[section]
\newtheorem{proposition}[theorem]{Proposition}
\newtheorem{lemma}[theorem]{Lemma}
\newtheorem{corollary}[theorem]{Corollary}
\theoremstyle{definition}\newtheorem{definition}[theorem]{Definition}
\theoremstyle{remark}\newtheorem{remark}[theorem]{Remark}
\newcommand{\OO}{\mathcal O}
\newcommand{\PP}{\mathbf P}
\newcommand{\FF}{\mathbf F}
\newcommand{\Pic}{\operatorname{Pic}}
\newcommand{\NS}{\operatorname{NS}}
\newcommand{\End}{\operatorname{End}}
\newcommand{\Hom}{\operatorname{Hom}}
\newcommand{\Nm}{\operatorname{Nm}}
\newcommand{\Tr}{\operatorname{Tr}}
\newcommand{\ann}{\operatorname{ann}}
\newcommand{\rk}{\operatorname{rank}}
\newcommand{\diag}{\operatorname{diag}}
\newcommand{\workbenchurl}{https://github.com/KokunoYumeto/schemes-and-moduli-workbench}
\newcommand{\recordurl}{https://github.com/KokunoYumeto/mathematics-workbenches}
\title{Grothendieck's generic Frobenius-rank question on mixed abelian surfaces}
\author{}
\date{23 November 2026}
\begin{document}
\maketitle
\begin{abstract}
Grothendieck's crystalline-cohomology notes ask how the generic stable Frobenius rank of hypersurface sections depends on their degree. For every abelian surface of $2$-rank one over an algebraically closed field of characteristic two and every ample bundle of positive ordinary polarization index, we prove that all sufficiently large complete systems $|M^{3e}|$ have generic curve defect two. No finite-field model or ordinary-index-zero ample class is required. We also determine exact ranks in every degree when the ordinary index is zero or is no larger than the local index, prove explicit bounds in the other cases, and reduce them by a maximal actual line-bundle root. For a pair $(A,M)$ over $\FF_q$, with $q=2^f$ and ordinary index $a>0$, the sufficient range is $e>1+4q^{\lceil2a/f\rceil}/2^a$. The unrestricted eventual theorem follows by spreading and openness of stable Frobenius rank. Every mixed surface with a finite-field model admits an ordinary-index-zero ample class, whereas a geometric-generic quotient with $\NS(A)=\mathbf Z[L]$ and indices $(1,0)$ admits none. For the remaining smaller degrees, a trace-zero first-variation calculation identifies a sufficient determinant test with a restricted surface cup product and a residual-pencil boundary, including its full twisted contribution. Uniform determinant nonvanishing remains open. The inputs include integral Dieudonn\'e duality, finite isotropic descent, principal cubic Cartier theory, nodal norms, Prym polarizations and Deuring--Shafarevich. The mathematics and text were produced and checked by GPT-6 Astra (ChatGPT Pro mode); the main rank and conditional determinant theorems received separate same-model reconstructions and full comparisons, and no human has checked the mathematics.
\end{abstract}

\section{The question and the two polarization indices}

This section defines the two polarization indices and states the eventual and exact-degree rank theorems and conditional determinant criterion.

Stable Frobenius rank measures the part of coherent cohomology on which repeated Frobenius remains invertible. Grothendieck uses it to study the unit-root contribution of hypersurface sections. In \S8.5.8 of his crystalline notes he considers its maximal generic value and eventual positivity; the continuation has the marginal comment ``Or cela est \'etabli au n\textsuperscript{o}~7'' \citep[cote~4, archive sheets~7--8, PDF images~8--9]{Grothendieck}. We determine eventual sharp ranks for every positive-index polarization on a surface of $2$-rank one, determine exact ranks in two large classes in every degree, and reduce the remaining smaller degrees to a precise primitive problem. The distinction matters because the generic defect depends on the integral polarization, not only on the slopes of the ambient surface.

Let $k$ be algebraically closed of characteristic two and let $A/k$ be an abelian surface with $2$-rank $f_A=1$. We call such a surface mixed. For any abelian variety $B$, its $2$-rank is defined by $|B[2](k)|=2^{f_B}$; for a smooth curve $C$, put $f_C=f_{\operatorname{Jac}(C)}$. Let $M$ be ample and put $D=\chi(M)=M^2/2$. Its polarization is
\[
 \lambda_M:A\longrightarrow A^\vee=\Pic^0(A),\qquad
 x\longmapsto t_x^*M\otimes M^{-1}.
\]
The geometric Tate modules $T_2(A)$ and $T_2(A^\vee)$ have rank one over $\mathbf Z_2$. Define
\[
 a=a(M)=\operatorname{length}_{\mathbf Z_2}\operatorname{coker}T_2(\lambda_M),
 \qquad b=b(M)=v_2(D)-a(M).
\]
We prove that $b\ge0$. These two integers are the elementary-divisor valuations of the ordinary and local-local integral polarization blocks. They need not be equal, and no product decomposition of $A$ is assumed.

For a smooth ample curve $C\subset A$, set
\[
 Q_C=\operatorname{coker}\bigl(H^1(A,\OO_A)\longrightarrow H^1(C,\OO_C)\bigr),
 \qquad r(C)=\dim_k\bigcap_{j\ge0}\operatorname{im}F^j|_{Q_C}.
\]
Frobenius is semilinear. The map defining $Q_C$ is injective. For $C\in|M^n|$, the elementary cohomology argument in Appendix~\ref{app:cohomology} gives
\begin{equation}\label{eq:dimensions}
 g(C)=Dn^2+1,\qquad \dim Q_C=Dn^2-1,\qquad r(C)=f_C-1.
\end{equation}
The generic stable rank is the constant maximal value on a nonempty open of the smooth-member locus. Assertions over a perfect field mean these geometric generic values after algebraic closure.

\begin{theorem}[Exact ranges]\label{thm:p10exact}
For each $d\ge1$, let $r_d$ be the generic relative stable rank in the complete system $|M^{3d}|$. Then:
\begin{enumerate}
\item If $a=0$, one has $r_d=9Dd^2-1$ for odd $d$ and $r_d=9Dd^2-2$ for even $d$.
\item If $0<a\le b$, one has $r_d=9Dd^2-2$ for every $d\ge1$.
\item If $a=1$ and $d$ is odd, one has $r_d=9Dd^2-2$, with no condition on $b$.
\end{enumerate}
The coherent relative dimension is $9Dd^2-1$ and the genus is $9Dd^2+1$. Thus the full curve defect $g-f$ is one in the first odd-degree case, and two in the other displayed cases.
\end{theorem}

\begin{theorem}[Integral root and bounds]\label{thm:p10bounds}
Put $c=\min(a,b)$. There is an ample bundle $H$ on the same surface with
\[
 M\simeq H^{2^c},\qquad a(H)=a-c,\qquad b(H)=b-c,
 \qquad \chi(H)=D/2^{2c}.
\]
The integer $c$ is maximal for existence of such an actual root in $\Pic(A)$. If $a>0$, then
\[
 \begin{array}{ll}
 9Dd^2-2^a\le r_d\le9Dd^2-2,&d\text{ odd},\\
 9Dd^2-2^{a+1}\le r_d\le9Dd^2-2,&d\text{ even}.
 \end{array}
\]
If $a>b>0$, the stronger uniform lower bound is
\[
                   r_d\ge9Dd^2-2^{a-b+1}\qquad(d\ge1).
\]
Consequently the unclassified mixed cases reduce to ample $H$ with $a(H)>0$ and $b(H)=0$, in the complete systems $|H^{3e}|$.
\end{theorem}

\begin{theorem}[Eventual sharp rank for every mixed pair]\label{thm:eventual-all-mixed}
Let $A/k$ be a mixed abelian surface over any algebraically closed field of characteristic two, and let $M$ be ample with $a(M)>0$. Put $D=\chi(M)$. There is an integer $e_0=e_0(A,M)$ such that, for every integer $e\ge e_0$, the generic smooth curve $C\in|M^{3e}|$ satisfies
\[
 f_C=9De^2-1,\qquad r(C)=9De^2-2,\qquad g(C)-f_C=2.
\]
No finite-field model or ordinary-index-zero ample class on $A$ is assumed. The threshold may depend on the pair.
\end{theorem}

\begin{theorem}[An ample difference]\label{thm:ample-difference}
Suppose $a(M)>0$ and that, on the same surface after algebraic closure, there is an ample line bundle $N$ with $a(N)=0$. For every $e\ge1$ such that $M^e\otimes N^{-1}$ is ample, the generic smooth curve in $|M^{3e}|$ has defect two and
\[
                         r_e=9\chi(M)e^2-2.
\]
In particular this holds for every sufficiently large $e$. An explicit sufficient bound is
\[
 e>\frac{M\cdot N+\sqrt{(M\cdot N)^2-M^2N^2}}{M^2};
\]
the stronger rational bound $e>(M\cdot N)/\chi(M)$ also suffices. Theorem~\ref{thm:no-ordinary-index-zero-ample-class} shows that such an $N$ need not exist on an arbitrary mixed surface.
\end{theorem}

\begin{theorem}[Finite-field models]\label{thm:finite-field}
Suppose $A/k$ is obtained by scalar extension from a mixed abelian surface over a finite field. Then $A$ admits an ample line bundle $N$ with $a(N)=0$. Consequently every ample $M$ with $a(M)>0$ has
\[
                         r_e=9\chi(M)e^2-2
\]
for all sufficiently large integers $e$, equivalently generic curve defect two in $|M^{3e}|$.

If the pair $(A,M)$ is defined over $\FF_q$, where $q=2^f$, put $a=a(M)>0$ and $m=\lceil2a/f\rceil$. The explicit sufficient range is
\begin{equation}\label{eq:finite-threshold}
                         e>1+\frac{4q^m}{2^a}.
\end{equation}
No optimality is asserted. The existence assertion does not require $\chi(N)$ to be odd, and the explicit bound requires descent of the pair, not merely of $A$.
\end{theorem}

\begin{theorem}[Failure of index-zero existence]\label{thm:no-ordinary-index-zero-ample-class}
There is a mixed abelian surface $A$ over an algebraically closed field of characteristic two and an ample line bundle $L$ such that
\[
 \NS(A)=\mathbf Z[L],\qquad \chi(L)=2,\qquad a(L)=1,\qquad b(L)=0.
\]
Every ample line bundle $N$ has class $t[L]$ for an integer $t>0$, and
\[
                  a(N)=1+v_2(t),\qquad b(N)=v_2(t).
\]
In particular $A$ has no ample line bundle of ordinary index zero and admits no finite-field model. This obstructs the hypothesis of Theorem~\ref{thm:ample-difference}, but Theorem~\ref{thm:eventual-all-mixed} still gives eventual defect two on this surface. It does not decide its unclassified smaller degrees or a determinant-nonvanishing assertion.
\end{theorem}

For a small numerical example, take $D=4$ and $a=b=1$. The root has Euler characteristic one, and the first system has genus $37$, full $2$-rank $35$, and relative stable rank $34$ in a $35$-dimensional space. By contrast, if $a=1,b=0,D=2$, the first system has genus $19$ and full $2$-rank $17$; in the next system the present bounds give relative rank between $68$ and $70$. The bounds alone do not choose an endpoint in that second example.

The integral inputs are classical. The splitting of finite group schemes is \citet[Lemma~4.44 and Proposition~4.45, p.~67]{EGM}; the integral Dieudonn\'e lattice, duality and Riemann form are in \citet[Theorem~3.14 and Corollaries~3.15--3.16, pp.~87--88; Propositions~3.18, 3.22 and 3.24, pp.~89, 99--101]{Oda1969}. Finite isotropic descent and the realization of symmetric homomorphisms by line bundles hold also in characteristic two \citep[Corollary~8.11, p.~117; Proposition~11.2 and Corollary~11.3, pp.~159--160]{EGM}. The curve-cover input is the unramified Deuring--Shafarevich formula, stated by \citet[equation~(1.8.2), p.~36]{Crew}. We give its short sheaf-filtration proof in the case used here. The integral Prym identity is due to \citet[\S3, equations~(3.2)--(3.3), pp.~346--347]{Debarre}. The characteristic-two cubic-principal argument is reproduced in the appendices so that the present rank deductions can be checked independently of the companion treatment.

The proof has two different geometric steps. An index-zero polarization has an \emph{\'etale principal norm source}, even when its Euler characteristic is even. Cubic principal curves on that source give sharp odd-degree nodal norms, and transverse unions give the even degrees. Integral Prym torsion supplies the matching obstruction. A common power of two in the two polarization blocks is an actual power of a line bundle, which transfers entire complete systems without restricting their parameter spaces. When the ordinary index is larger, an \emph{\'etale quotient} instead produces test curves whose defects are multiplied by a power of two; those give bounds for every degree. The remaining fixed-degree problem concerns the cases outside the exact ranges; Theorem~\ref{thm:eventual-all-mixed} settles all sufficiently large degrees for each pair.

Theorem~\ref{thm:ample-difference} is a further consequence of the ordinary-index-zero result and nodal smoothing: the two ample summands each have a sharp defect-one cubic system. It narrows the remaining rank problem on surfaces admitting an index-zero class, without proving the nonvanishing of the particular cup-product determinant below. The numerical ampleness bound follows directly from abelian Riemann--Roch and the criterion for an effective nondegenerate line bundle to be ample \citep[Proposition~2.22, p.~25; Theorem~9.11, pp.~131--133]{EGM}. Theorem~\ref{thm:finite-field} removes the existence hypothesis on every mixed surface with a finite-field model. Its proof combines Frobenius--Verschiebung duality, the nilpotence bound for a finite connected coordinate algebra, the actual-root theorem and the theorem of the cube. It establishes eventual sharpness without asserting the missing determinant condition or settling every smaller degree. The finite-field construction and its explicit bound are proved in Section~\ref{sec:finite-field}. Theorem~\ref{thm:eventual-all-mixed} removes the finite-field hypothesis from the eventual rank conclusion. Section~\ref{sec:spreading} proves it from the finite-field theorem by standard spreading, relative cohomology and an explicit stable-Frobenius-minor argument \citep[Tags~01ZM, 0B91, 02O1 and 0D4E]{Stacks}. This is an existential conclusion for each pair; it does not provide a uniform threshold or prove the actual determinant condition. Theorem~\ref{thm:no-ordinary-index-zero-ample-class} shows that the finite-field existence result cannot extend to every mixed surface. Its input is the full symplectic monodromy theorem for the $2$-rank-one genus-two stratum of \citet[Theorem~4.5, author version pp.~14--15]{AchterPries}. The geometric-generic endomorphism argument is classical \citep[Corollary~4.2, p.~297]{Chai}; the new conclusion here is the exact ordinary index on every ample class of the quotient. Section~\ref{sec:no-zero-class} gives the complete construction and the elementary $\alpha_2$ theta-descent proof.

The results concern complete systems. If the embedding is given by $M^3$, restriction from ambient degree-$d$ equations is surjective for $d\ge\max\{1,\operatorname{reg}(\mathcal I_A)-1\}$, and in that range the same generic ranks hold for ambient hypersurfaces. No claim about a monodromy group or arbitrary ambient varieties is made; the remaining primitive degrees and the exact determinant question are stated below.

The remaining question admits a further conditional theorem. In its statement, a multiplicative $\mu _2$-linearization of a positive-ordinary-index bundle gives the idempotent weight operator $e$ on its sections. A section $s$ is called sharp when $sF$ is bijective on the annihilator of $\langle s,es\rangle$ in $H^2(H^{-1})$.

\begin{theorem}[Conditional determinant reduction]\label{thm:intro-reduction}
Let $B/k$ be a mixed abelian surface, let $H=M^{3e}$ be ample with positive ordinary index, and let $\pi:\widetilde B\to B$ be a connected \'etale double isogeny. Suppose a sharp section of $H$ exists. At its pullback, restricted coherent Frobenius has first and stable corank two. Its kernel-to-cokernel first variation factors through
$H^0(B,H)/\langle s,es\rangle$ and is trace zero in the compatible torsor-cup bases. If every member is singular, its image is zero or one rank-one nilpotent matrix line. An invertible member produces a sharp member of the complete system $|\pi^*H|$.

For a general sharp pair $s=s_0+s_1$ with disjoint zero schemes, factor relative Frobenius as $B\xrightarrow{\rho}C\xrightarrow{\nu}T=B^{(2)}$, with kernels $\mu_2$ and $\alpha_2$. Descend $H=\rho^*J$, put $L=H^{(2)}$, and let $P$ be the nontrivial geometric two-torsion line bundle on $T$. The Cartier-paired summands are $E_0=\nu_*J$ and $E_1=E_0P$. Let $K_s$ be the kernel of their contraction by $(s_0,s_1)$, and let $Q_s=K_s/P$, where $P$ maps by $z\mapsto(zs_1,zs_0)$ using a fixed $P^2\simeq\OO_T$. Put $\mathcal L=LP=\det Q_s$ and $U_s=H^1(T,Q_s^\vee)$; this space has dimension two. Define
\[
 \Delta_s=\{s_0\wedge t_1+t_0\wedge s_1:t_i\in H^0(T,E_i)\}
       \subset H^0(T,\mathcal L),
\]
\[
 \mu_{\Delta}:\operatorname{Sym}^2 U_s\longrightarrow\Delta_s^*,
 \qquad xy\longmapsto
       \bigl[v\mapsto\operatorname{Tr}_T(v\,\wedge_*(x\cup y))\bigr].
\]
Here the wedge pairing takes values in $\det Q_s^\vee=\mathcal L^{-1}$, and an invariant volume form fixes the trace. An invertible actual first variation exists if and only if
\[
 \{x\in U_s:\mu_{\Delta}(xy)=0\text{ for every }y\in U_s\}=0.
\]
The determinant-direction map has kernel $H^0(T,K_sP)$ and cokernel $H^1(T,K_sP)$. No assertion of uniform vanishing of this common radical is part of the theorem.
\end{theorem}

In words, the missing rank increment is reduced to a two-dimensional obstruction with a rigid trace identity. The curve construction specifies that obstruction, including the twisting which a bare dimension count would lose. Theorem~\ref{thm:boundary} gives its residual-pencil form with the full two-step boundary kernel; its three-dimensional product-space assertion assumes that the determinant divisor is integral. Proposition~\ref{prop:auto-directions} also computes an unavoidable automorphism contribution to the actual direction codimension; in positive local index that codimension is at least two. This does not prove determinant nonvanishing. Sections~\ref{sec:toral}--\ref{sec:open} prove the theorem.

The polynomial in the cubic-principal appendix is already explicit in \citet[equation~(3.9), author version p.~14]{GrusonSam}; their geometric treatment excludes characteristics two and three. Its characteristic-two Cartier application is proved here. For recent neighboring work, \citet[\S1.3 and acknowledgements]{Zou} studies Frobenius splitting and crystalline slope zero for higher-dimensional primitive symplectic varieties and records AI assistance in those arguments. Those results concern a different cohomological setting and are not inputs to our curve-rank formulas. A bounded search of arXiv, journals, author pages, ECCC, Hexagon, and recent AI mathematics did not identify these exact mixed-polarization formulas or the present restricted determinant criterion.

\section{The integral polarization and its maximal root}

We first justify that the two indices control an integral homomorphism, not merely a rational slope decomposition.

\begin{lemma}\label{lem:p10blocks}
The ordinary and local-local blocks of the integral polarization pairing have elementary-divisor valuations $(a,a)$ and $(b,b)$, respectively. The finite group scheme $\ker\lambda_M$ is killed by $D$. There is an ample bundle $N$ on $A^\vee$ such that
\[
 \lambda_M^*N\simeq M^D,\qquad
 \lambda_N=D\lambda_M^{-1},\qquad
 \chi(N)=D,\qquad a(N)=b,\quad b(N)=a.
\]
\end{lemma}
\begin{proof}
Over a perfect field the connected--\'etale sequence of a finite commutative group scheme splits functorially: its reduced subgroup is its \'etale factor. A homomorphism between a local and an \'etale group is zero. These statements and their proofs are \citet[Lemma~4.44 and Proposition~4.45, p.~67]{EGM}. Apply the splitting at every level of $A[2^\infty]$, and apply it also to the Cartier dual of the connected factor. The compatible decompositions give
\[
 A[2^\infty]=G_{\mathrm{et}}\times G_{\mathrm{mult}}\times G_{\mathrm{ll}},
 \quad G_{\mathrm{et}}\simeq\mathbf Q_2/\mathbf Z_2,
 \quad G_{\mathrm{mult}}\simeq\mu_{2^\infty}.
\]
The last factor and its dual are connected; it has height two and dimension one. All cross-factor homomorphisms vanish, including the multiplicative/local-local ones by Cartier duality. The same decomposition holds on the dual surface.

Use the covariant Dieudonn\'e lattice, the $W(k)$-linear dual of the contravariant convention of \citet{Oda1969}. It is free, exactness identifies finite-group orders with lattice lengths, and Cartier duality is integral lattice duality at the cited locations. A polarization gives a skew-symmetric integral bilinear form \citep[Propositions~1.10 and 1.12, pp.~75--76; Proposition~3.24(ii), pp.~100--101]{Oda1969}. Since $W(k)$ has no $2$-torsion, the form is alternating. The ordinary and local-local summands are orthogonal by cross-factor vanishing.

The ordinary block pairs its two rank-one factors, and integral duality makes their two indices equal to the geometric Tate index $a$. On the rank-two local-local lattice the matrix is
\[
 \begin{pmatrix}0&u\\-u&0\end{pmatrix},\qquad u\in W(k)\setminus\{0\}.
\]
Writing $b=v_2(u)$ gives two equal elementary-divisor valuations. Since $\deg\lambda_M=D^2$, the length formula gives $2v_2(D)=2a+2b$. This also proves $b\ge0$. The two-primary kernel is killed by $2^{\max(a,b)}$, hence by the two-primary part of $D$.

For an odd prime $\ell$, the kernel is constant with a perfect alternating commutator pairing. An element of maximal order $\ell^e$ has a partner whose pairing has that order. Pairing a relation with both generators shows that they generate a subgroup of order $\ell^{2e}$; hence $e\le v_\ell(D)$. Thus the entire kernel is killed by $D$.

On $\ker\lambda_M$, the commutator pairing of $M^D$ is the $D$th power of that of $M$, hence trivial. Finite isotropic descent \citep[Corollary~8.11 and its proof, p.~117]{EGM} gives $N$ with the asserted pullback. It is ample because its finite pullback is ample. Pullback polarizations give
\[
 \lambda_M\lambda_N\lambda_M=D\lambda_M.
\]
Cancellation gives the integral identity $\lambda_N=D\lambda_M^{-1}$. Riemann--Roch gives $\chi(N)=D^3/D^2=D$. On each block the inverse matrix multiplied by $D$ replaces its valuation by $v_2(D)$ minus that valuation, proving the exchanged indices.
\end{proof}

\begin{lemma}\label{lem:p10root}
The root assertion of Theorem~\ref{thm:p10bounds} holds.
\end{lemma}
\begin{proof}
Put $n=2^{\min(a,b)}$. The two integral blocks show that $\lambda_M$ annihilates the entire group scheme $A[n]$. Since $[n]:A\to A$ is its faithfully flat quotient, there is an integral homomorphism $\varphi:A\to A^\vee$ with $\lambda_M=n\varphi$. It is symmetric: dualizing gives $n(\varphi-\varphi^\vee)=0$, and $\Hom(A,A^\vee)$ has no integer torsion. Indeed a torsion homomorphism factors through a finite affine group scheme, and every morphism from the proper connected reduced variety $A$ to a finite affine scheme is constant.

Every symmetric homomorphism over an algebraically closed field is induced by a line bundle, including in characteristic two. The complete finite-isotropic-descent proof is \citet[Proposition~11.2 and Corollary~11.3, pp.~159--160]{EGM}. Choose $H_0$ with $\lambda_{H_0}=\varphi$. Then $M\otimes H_0^{-n}\in\Pic^0(A)$. Multiplication by $n$ on the dual abelian variety is surjective on algebraically closed points, so a twist of $H_0$ by an algebraically trivial bundle gives $H^n\simeq M$. A line bundle whose positive tensor power is ample is ample. The polarization identity subtracts $v_2(n)$ from both indices, and $\chi(H^n)=n^2\chi(H)$ gives the claimed Euler characteristic. Conversely any root of order $2^s$ forces both block valuations to be at least $s$, proving maximality.
\end{proof}

\section{Principal norm sources and the connected Prym}

This section constructs the principal norm source and proves the integral Prym obstruction needed for the exact ranks.

\begin{lemma}\label{lem:p10source}
The equality $a(M)=0$ holds if and only if there exist a principally polarized surface $(B,\Theta)$ and a finite \'etale isogeny $\phi:B\to A$ of degree $D$ with
\[
 \phi^*M\equiv D\Theta,\qquad \Nm_\phi\OO_B(\Theta)\equiv M.
\]
The surface $B$ has $2$-rank one.
\end{lemma}
\begin{proof}
Suppose $a=0$ and use $N$ from Lemma~\ref{lem:p10blocks}. Its local-local kernel is trivial, while its ordinary two-primary kernel is the product of a multiplicative group and its \'etale dual, each of order $2^b$. Each factor is isotropic: the restriction of the commutator pairing is a map between a local group and an \'etale dual. Perfection makes the multiplicative factor maximal isotropic. Combine it with an isotropic subgroup of half-order in each odd-primary symplectic factor. The resulting subgroup $H\subset\ker\lambda_N$ has order $D$ and \'etale Cartier dual.

Finite descent gives $\psi:A^\vee\to C=A^\vee/H$ and a principal polarization $\Theta_C$ with $\psi^*\Theta_C\simeq N$ \citep[Proposition~8.19 and proof, p.~119]{EGM}. Put $B=C^\vee$ with its dual principal polarization and $\phi=\psi^\vee$. Isogeny duality identifies its kernel with $H^D$ \citep[Theorem~7.5, p.~100]{EGM}, so it is \'etale of degree $D$. The identity $\psi\lambda_M\phi=D\lambda_C^{-1}$ gives the first numerical equality. Pullback of the norm of $\Theta$ is the sum of its $D$ translates, numerically $D\Theta$; injectivity of pullback on the torsion-free N\'eron--Severi group gives the second. Isogenies preserve $2$-rank.

Conversely an \'etale isogeny induces an isomorphism on local-local $2$-divisible factors. Comparing those blocks in $\phi^*M\equiv D\Theta$ gives $b(M)=v_2(D)$, hence $a(M)=0$.
\end{proof}

\begin{lemma}[Integral Prym obstruction]\label{lem:p10prym}
For every smooth $C\in|M^m|$, its full defect is at least one. If $a(M)+v_2(m)>0$, it is at least two.
\end{lemma}
\begin{proof}
Restriction $j:A^\vee\to J_C$ is a closed immersion. For nontrivial $P\in\Pic^0(A)(k)$, one has $H^0(A,P)=0$ and $H^1(A,P(-C))=0$ by ample vanishing and duality, so $P|_C$ cannot be trivial. Its tangent kernel also vanishes by $H^1(A,\OO_A(-C))=0$. Thus the finite kernel is \'etale with only its identity point. Let $q:J_C\to A$ be the dual map under the Jacobian polarization. Its kernel $P_C$ is smooth and connected: after quotienting by the reduced identity component, the remaining map is an isogeny, whose dual kernel embeds in $\ker j=0$. Therefore $\dim P_C=g(C)-2$ and $f_C=f_{P_C}+1$, using isogeny invariance and Poincar\'e reducibility. This proves the first bound.

The exact induced-polarization formula is
\[
 qj=j^*\lambda_{J_C}=mD\lambda_M^{-1}=m\lambda_N.
\]
Indeed the norm identity in \citet[\S3, equations~(3.2)--(3.3), pp.~346--347]{Debarre} gives this after cancellation in the rational homomorphism group, and both sides are integral. The local-local block of $\lambda_N$ has index $a(M)$. If $a(M)+v_2(m)>0$, the homomorphism $m\lambda_N$ kills the full level-one local-local group $G_{\mathrm{ll}}[2]$ on $A^\vee$. It therefore embeds into $P_C[2]$ through $j$. After Frobenius twist, functorial Verschiebung embeds the connected order-two kernel of $V$ on this factor into $\ker V_{P_C}$. For an abelian variety $P$, the connected part of $\ker V_P$ has order $2^{\dim P-f_P}$; this follows from the Frobenius--Verschiebung identities and the torsion description \citep[Propositions~5.20--5.22, pp.~80--81]{EGM}. Thus $\dim P_C-f_{P_C}\ge1$, proving the second bound.
\end{proof}

\begin{proof}[Proof of Theorem~\ref{thm:p10exact}, parts~1 and~2]
Suppose first $a=0$ and $d$ is odd. Apply Lemma~\ref{lem:p10source} to $M^d$, whose Euler characteristic is $Dd^2$ and whose ordinary index remains zero. Proposition~\ref{prop:seed} supplies a smooth cubic principal curve of genus ten and $2$-rank nine on that source, and Lemma~\ref{lem:norm-nodes} supplies a nonempty open of curves whose norm images are integral nodal. The two opens meet. The resulting image has $9(Dd^2-1)$ nodes, genus $9Dd^2+1$, and defect one by Lemma~\ref{lem:nodal}. Its numerical class is $M^{3d}$; Lemma~\ref{lem:exact-system} translates it into that exact system and proves $M^{3d}$ very ample. Stable-rank openness gives a smooth member of defect at most one; Lemma~\ref{lem:p10prym} forces equality.

For even $d$, take sharp smooth members of $|M^3|$ and $|M^{3(d-1)}|$ meeting transversely. Very ampleness and the open rank condition permit this choice. Their connected nodal union has defect two. Openness supplies a smooth member of defect at most two, and the Prym bound gives equality. Subtracting the ambient $2$-rank one proves part~1.

If $0<a\le b$, Lemma~\ref{lem:p10root} gives $M=H^{2^a}$ with $a(H)=0$. The complete systems are identical to $|H^{3d'}|$, where $d'=2^ad$ is even. Part~1 gives
\[
 r_d=9\chi(H)(d')^2-2=9(D/2^{2a})(2^ad)^2-2=9Dd^2-2.
\]
This proves part~2 without any claim about genericity in a restricted pullback subspace.
\end{proof}

\section{\'Etale quotient curves and the bounds}

This section constructs quotient curves at positive ordinary index and uses their defects to prove the remaining rank bounds.

\begin{lemma}\label{lem:p10quotient}
If $a>0$ and $q=2^a$, there are a cyclic \'etale isogeny $\pi:A\to B$ of degree $q$ and an ample bundle $J$ on $B$ with
\[
 M\simeq\pi^*J,\qquad a(J)=0,\qquad \chi(J)=D/q,
 \qquad b(J)=b(M).
\]
Every smooth $C\in|J^m|$ has smooth connected pullback $Y$, and
\[
                       g(Y)-f_Y=q\bigl(g(C)-f_C\bigr).
\]
\end{lemma}
\begin{proof}
The ordinary \'etale factor of $\ker\lambda_M$ is the constant cyclic group $K=\mathbf Z/q\mathbf Z$. Quotient by it. To descend the actual line bundle, choose a generator $x$ and an isomorphism $t_x^*M\simeq M$. Its $q$th iterate is multiplication by a scalar; rescale the isomorphism by a $q$th root of the inverse scalar. Its powers now give a genuine $K$-linearization. Faithfully flat descent gives $J$, which is ample because its pullback is ample. This also proves the isotropic descent assertion here without division by the group order.

In $\lambda_M=\pi^\vee\lambda_J\pi$, the map $T_2(\pi)$ has cokernel length $a$, while $T_2(\pi^\vee)$ is an isomorphism because its kernel is multiplicative. Thus $a(J)=0$. The local-local factors are unchanged by an \'etale isogeny, giving $b(J)=b(M)$; Euler characteristics give $\chi(M)=q\chi(J)$.

The pullback is smooth by \'etaleness and connected because it is an ample divisor: the divisor sequence and antiample $H^1$ vanishing give $H^0(Y,\OO_Y)=k$. It is therefore a connected cyclic \'etale $2$-power cover of $C$. The Deuring--Shafarevich formula gives $f_Y-1=q(f_C-1)$ \citep[Theorem~1.5, Corollaries~1.7--1.8 and equation~(1.8.2), pp.~34--36]{Crew}. Together with \'etale Hurwitz this gives the defect formula.

Here is its sheaf proof in precisely this case. The locally constant sheaf $\pi_*\mathbf F_2$ has the augmentation filtration of the regular module $\mathbf F_2[K]=\mathbf F_2[t]/(t^q)$, with $q$ trivial one-dimensional quotients. Additivity of Euler characteristic therefore gives $\chi_{\mathrm{et}}(Y,\mathbf F_2)=q\chi_{\mathrm{et}}(C,\mathbf F_2)$. Artin--Schreier theory gives $H^i_{\mathrm{et}}(C,\mathbf F_2)=0$ for $i>1$ and $\dim H^1_{\mathrm{et}}(C,\mathbf F_2)=f_C$, so its Euler characteristic is $1-f_C$. The same applies to $Y$. This is the argument of \citet[\S3, first alternative proof, p.~42; Lemma~1.9.2, p.~37]{Crew}, and yields the formula directly.
\end{proof}

\begin{proof}[Proof of Theorem~\ref{thm:p10bounds} and part~3 of Theorem~\ref{thm:p10exact}]
The root was proved in Lemma~\ref{lem:p10root}. For $a>0$, apply Lemma~\ref{lem:p10quotient} and take a sharp smooth member of $|J^{3d}|$ from the index-zero theorem. Its defect is one for odd $d$ and two for even $d$. Its smooth connected pullback in $|M^{3d}|$ has defect $2^a$ or $2^{a+1}$, respectively. Stable-rank openness gives the two lower bounds; Lemma~\ref{lem:p10prym} gives the common upper bound. When $a=1$ and $d$ is odd, the bounds coincide.

If $a>b>0$, use the exact root with exponent $2^b$. Its ordinary index is $a-b$, and the corresponding degree $2^bd$ is even. The even lower bound just proved gives
\[
 r_d\ge9\chi(H)(2^bd)^2-2^{a-b+1}
       =9Dd^2-2^{a-b+1}.
\]
Finally the root reduces $a>b$ to $(a(H),b(H))=(a-b,0)$. For $a\le b$ the exact ranges already apply. All identifications are of complete systems of isomorphic line bundles on the same surface.
\end{proof}

\section{An ample difference and eventual sharpness}
\label{sec:ample-difference}

This section proves the exact rank in an additional range by joining two curves of defect one. The construction uses an ordinary-index-zero class on the same surface and does not require nonvanishing of the restricted cup-product determinant considered below.

\begin{proof}[Proof of Theorem~\ref{thm:ample-difference}]
Work over an algebraic closure of the ground field and put
\[
                 H_e=M^e\otimes N^{-1}.
\]
Choose bases of the two rank-one geometric $\mathbf Z_2$ Tate modules. The homomorphisms induced by $\lambda_M$ and $\lambda_N$ are multiplication by scalars $c_M$ and $c_N$, with $2\mid c_M$ and $c_N$ a unit. Additivity of polarization homomorphisms gives the scalar
\[
                         e c_M-c_N
\]
for $\lambda_{H_e}$. It is a unit for every integer $e$. Thus $a(H_e)=0$ whenever $H_e$ is ample.

The ordinary-index-zero case of Theorem~\ref{thm:p10exact}, applied in degree one to $N$ and $H_e$, gives nonempty opens of smooth curves of defect one in $|N^3|$ and $|H_e^3|$. Choose such a curve $C_0$ in the first system. Since $H_e^3$ is very ample, the curves transverse to $C_0$ form a nonempty open in the second system. This open meets its sharp smooth open, so we can choose $C_1$ satisfying both conditions. The curves meet at
\[
                      q=9(N\cdot H_e)>0
\]
distinct points. Their union $C_*=C_0\cup C_1$ is a connected nodal divisor in $|M^{3e}|$.

The normalization sequence and its degree-one cohomology are
\[
 0\longrightarrow\OO_{C_*}\longrightarrow
 \OO_{C_0}\oplus\OO_{C_1}\longrightarrow k^q\longrightarrow0,
\]
\[
 0\longrightarrow k^{q-1}\longrightarrow H^1(C_*,\OO_{C_*})
 \longrightarrow H^1(C_0,\OO_{C_0})\oplus H^1(C_1,\OO_{C_1})
 \longrightarrow0.
\]
The image of $k^2\to k^q$ is the diagonal line. Both sequences commute with Frobenius, which is bijective on $k^{q-1}$. Exactness of the stable part for finite-dimensional semilinear spaces, proved in Appendix~\ref{app:cohomology}, therefore gives
\[
\begin{split}
 f(C_*)&=f(C_0)+f(C_1)+q-1,\\
 p_a(C_*)&=g(C_0)+g(C_1)+q-1.
\end{split}
\]
Here $f(C_*)$ denotes the stable coherent Frobenius dimension on the nodal curve. It follows that $p_a(C_*)-f(C_*)=2$.

The universal divisor over the complete system $|M^{3e}|$ is flat and projective. Near $C_*$ its fibres have $h^0(\OO)=1$: constants give the lower bound, and upper semicontinuity gives the upper bound. Flatness fixes the Euler characteristic, so $h^1(\OO)$ is constant there. Cohomology and base change give a vector bundle $\mathcal H$ of rank $g=p_a(C_*)$, with its relative Frobenius map $F_S^*\mathcal H\to\mathcal H$.

On every geometric fibre the stable rank is the rank of the $g$th iterate. Indeed the ranks of successive images decrease, and equality of two successive ranks makes the image Frobenius-stable over the perfect residue field. The condition that this iterate have rank at least $g-2$ is open, expressed by its minors. It holds at $C_*$ and hence on a nonempty open of the irreducible complete system. The smooth locus is another nonempty open, because $M^{3e}$ is very ample. Their intersection contains smooth curves of defect at most two.

The upper bound in Theorem~\ref{thm:p10bounds} gives defect at least two for the generic smooth member. Thus its defect is exactly two. Finally
\[
 g=9\chi(M)e^2+1,\qquad r_e=f-1=9\chi(M)e^2-2,
\]
as required. The equality is geometric, so it gives the stated generic rank over every perfect ground field.
\end{proof}

\begin{lemma}\label{lem:ample-positive-square}
Let $M$ be ample on an abelian surface. A line bundle $R$ with $R^2>0$ and $M\cdot R>0$ is ample. If $N$ is also ample, then
\[
                         (M\cdot N)^2\ge M^2N^2.
\]
\end{lemma}
\begin{proof}
The abelian Riemann--Roch identities, valid for every line bundle, are
\[
 \chi(R)=R^2/2,\qquad \deg\lambda_R=\chi(R)^2;
\]
the degree is zero when the polarization homomorphism is not finite \citep[Theorem~9.11 and proof, pp.~131--133]{EGM}. Since $M\cdot R>0$, the line bundle $R^{-1}$ has no nonzero section: a nonzero effective divisor has positive intersection with an ample class. The canonical bundle is trivial, so Serre duality gives $h^2(R)=0$. Therefore
\[
                     h^0(R)=\chi(R)+h^1(R)>0.
\]
The same Riemann--Roch identities make $\ker\lambda_R$ finite. An effective line bundle with finite polarization kernel is ample \citep[Proposition~2.22 and proof, p.~25; Lemma~2.21, pp.~24--25]{EGM}. This proves the first assertion.

For the second, set $x=M^2$, $y=M\cdot N$, $z=N^2$, and use additive divisor notation for $R=xN-yM$. Then $M\cdot R=0$ and $R^2=x(xz-y^2)$. If $xz-y^2>0$, Riemann--Roch and Serre duality imply
\[
                   h^0(R)+h^0(R^{-1})\ge\chi(R)>0.
\]
One of $R,R^{-1}$ would have a nonzero effective divisor of intersection zero with $M$, a contradiction. Hence $y^2\ge xz$.
\end{proof}

\begin{proof}[Proof of the numerical bound in Theorem~\ref{thm:ample-difference}]
Write $x=M^2=2\chi(M)$, $y=M\cdot N$, and $z=N^2$. These are positive and $y^2\ge xz$. If
\[
                     e>\frac{y+\sqrt{y^2-xz}}{x},
\]
the class $eM-N$ has positive square $e^2x-2ey+z$ and positive intersection $ex-y$ with $M$. Lemma~\ref{lem:ample-positive-square} makes it ample. Since $\sqrt{y^2-xz}<y$, the stronger rational inequality $e>2y/x=(M\cdot N)/\chi(M)$ also suffices. This proves an explicit eventual range.
\end{proof}

The hypothesis on $N$ is essential to this argument. It is not asserted that every mixed surface admits an ordinary-index-zero ample class. When such a class exists, Theorem~\ref{thm:ample-difference} settles all sufficiently large degrees for every positive-index polarization, including degrees outside the ordered-index ranges. It says nothing about the value of the first-variation determinant at a prescribed sharp source. The unresolved cases include the remaining small degrees and surfaces on which the ordinary Tate homomorphism of every line bundle is divisible by two.

\section{Finite-field models and an explicit eventual range}
\label{sec:finite-field}

This section constructs an ordinary-index-zero ample bundle on every mixed surface with a finite-field model, and bounds the ample difference needed for the sharp rank formula.

We use two elementary consequences of Frobenius and Verschiebung. If $G$ is a finite connected group scheme of order $h$ over $\FF_q$, its coordinate algebra is local with augmentation ideal $I$ and $I^h=0$. Indeed, every nonzero term in the sequence $I,I^2,\ldots$ strictly decreases dimension until it is zero: equality of two successive terms would imply their vanishing by Nakayama's lemma. Thus the $q^m$-power Frobenius kills $G$ whenever $q^m\ge h$. If the Cartier dual $G^D$ is also connected, the same argument on $G^D$, followed by Cartier duality, shows that the $m$th $q$-Verschiebung kills $G$ as well. The duality used here is \citet[Proposition~5.19 and proof, pp.~79--80]{EGM}. Over a perfect field the \'etale, multiplicative and local-local factors are canonical \citep[Lemma~4.44, Proposition~4.45 and the following paragraph, pp.~67--68]{EGM}, so these statements apply to their factors in an abelian variety's torsion.

\begin{proof}[Proof of Theorem~\ref{thm:finite-field}]
First suppose that the pair $(A,M)$ is defined over $\FF_q$, where $q=2^f$, and that $a=a(M)>0$. All line-bundle roots below are taken after algebraic closure. Let $\pi$ be the $q$-Frobenius of this model and let $v$ be its Verschiebung, so that
\[
                     \pi v=v\pi=[q].
\]
These identities and their iterates are \citet[Proposition~5.20 and \S5.21, p.~80]{EGM}. Put
\[
 m=\left\lceil\frac{2a}{f}\right\rceil,
 \qquad u=\pi^m+v^m,\qquad w=\pi^m-v^m.
\]
The local-local factor $G$ of $A[2^a]$ has order $2^{2a}$, and both $G$ and $G^D$ are connected. Since $q^m\ge2^{2a}$, the preceding coordinate-algebra argument proves
\begin{equation}\label{eq:finite-kills}
                        \pi^mG=v^mG=0,
                        \qquad uG=0.
\end{equation}
This argument takes place on the finite-field model. It does not identify an arbitrary scalar extension of a relative Frobenius with the absolute power map on all scalar coefficients.

Because $\lambda_M$ is defined over $\FF_q$, Frobenius--Verschiebung duality gives
\[
 \pi^\vee\lambda_M=\lambda_Mv,
 \qquad v^\vee\lambda_M=\lambda_M\pi,
 \qquad u^\vee\lambda_M=\lambda_Mu;
\]
see \citet[Proposition~7.34 and proof, p.~111]{EGM}. On the rank-one \'etale Tate module, $\pi$ acts by a unit $z\in\mathbf Z_2^\times$ and $v$ by $q/z$. Thus $u$ acts by the unit
\[
                            z^m+q^m/z^m.
\]
Consider the actual line bundle
\begin{equation}\label{eq:finite-L}
                         L=u^*M\otimes M^{2^a}.
\end{equation}
A pullback of an ample line bundle is nef, and its tensor product with an ample bundle is ample. Hence $L$ is ample. Its polarization is
\[
 \lambda_L=u^\vee\lambda_Mu+2^a\lambda_M
           =\lambda_M(u^2+[2^a]).
\]
The ordinary Tate scalar has valuation $a$, because
$(z^m+q^m/z^m)^2+2^a$ is odd. Consequently $a(L)=a$.

The homomorphism $\lambda_L$ kills the whole group scheme $A[2^a]$. On its ordinary factor, $\lambda_M$ already kills this level by Lemma~\ref{lem:p10blocks}; on its local-local factor the first summand vanishes by \eqref{eq:finite-kills}. The second summand $2^a\lambda_M$ vanishes at this level on every factor. The two-block criterion therefore gives $b(L)\ge a$. The actual-root assertion of Theorem~\ref{thm:p10bounds} supplies an ample bundle $N$ with
\begin{equation}\label{eq:finite-root}
                         N^{2^a}\simeq L,
                         \qquad a(N)=0.
\end{equation}

We next prove the explicit estimate. Numerical equivalence of line bundles is written additively in $\NS(A)\otimes\mathbf Q$. Finite-field Frobenius gives $(\pi^m)^*[M]=q^m[M]$. Applying $(v^m)^*$ and using $\pi^mv^m=[q^m]$ gives $(v^m)^*[M]=q^m[M]$, since $[q^m]^*[M]=q^{2m}[M]$. The multiplication formula is \citet[Corollary~2.12 and proof, p.~21]{EGM}. The numerical parallelogram identity, a consequence of the theorem of the cube, now gives
\begin{equation}\label{eq:finite-parallelogram}
                  u^*[M]+w^*[M]=4q^m[M].
\end{equation}
Here $[-1]^*[M]=[M]$ numerically; the requisite pullback identity and its proof are \citet[Theorem~2.7 and Corollary~2.8, pp.~19--20]{EGM}. Combining \eqref{eq:finite-L}--\eqref{eq:finite-parallelogram} yields
\begin{equation}\label{eq:finite-nef}
 \left(1+\frac{4q^m}{2^a}\right)[M]-[N]
                  =\frac{1}{2^a}w^*[M],
\end{equation}
which is nef. For every integer $e>1+4q^m/2^a$, the class of $M^e\otimes N^{-1}$ is the sum of a positive multiple of the ample class $[M]$ and the nef class in \eqref{eq:finite-nef}. It is therefore ample. Theorem~\ref{thm:ample-difference} proves the stated defect and relative rank.

Finally suppose only that $A/k$ is the scalar extension of an abelian surface over a finite field. Choose an ample line bundle on that projective finite-field model, enlarging the finite field if necessary, and call it $P$. If $a(P)=0$, it is already the desired bundle. Otherwise the construction just given, with $P$ in place of $M$, produces an ample ordinary-index-zero root over the algebraic closure of that finite field. The bundle and its root isomorphism descend to a finite extension by finite presentation; base changing them to $k$ proves the first assertion. For an arbitrary ample $M$ on $A/k$ with $a(M)>0$, the previously constructed $N$ makes $M^e\otimes N^{-1}$ ample for all sufficiently large $e$. Theorem~\ref{thm:ample-difference} again gives the asserted eventual rank. The numerical bound uses a finite-field model of the pair $(A,M)$ and is asserted only under that additional hypothesis.
\end{proof}

The argument makes the ordinary polarization block invertible. It does not assert that $\chi(N)$ is odd: the local-local index may remain positive. Its rank conclusion is eventual, and the numerical threshold is only sufficient. Neither the unclassified smaller powers nor the uniform first-variation determinant are settled by this construction.

\section{Eventual sharp rank over arbitrary algebraically closed fields}
\label{sec:spreading}

This section transfers the finite-field rank theorem to every mixed pair by spreading the pair and using an open Frobenius-rank condition.

We use standard limit and cohomology results in their scheme-theoretic forms. Schemes of finite presentation, their morphisms and line bundles descend along filtered limits; smoothness, properness and relative ampleness hold after shrinking a finite-type model \citep[Tags~01ZM, 0B8W, 0C0C, 081F and 09MT]{Stacks}. For a proper finitely presented family and a finitely presented sheaf flat over the base, derived cohomology is perfect and commutes with arbitrary base change \citep[Tag~0B91]{Stacks}. Euler characteristic is locally constant \citep[Tag~0B9T]{Stacks}. Relative Serre vanishing supplies one threshold for all sufficiently large powers, and the resulting degree-zero direct image is locally free with arbitrary base change \citep[Tags~02O1 and 0D4E]{Stacks}.

\begin{lemma}\label{lem:stable-rank-open}
Let $\mathcal H$ be a vector bundle of rank $g$ on a scheme $T$ of characteristic two, with a map $\Phi:F_T^*\mathcal H\to\mathcal H$. The locus where the stable Frobenius rank on a geometric fibre is at least $q$ is open. It is the rank-at-least-$q$ locus of
\[
 \Phi^{[g]}:(F_T^g)^*\mathcal H\longrightarrow\mathcal H.
\]
In particular, $2$-rank at least $q$ is an open condition in a smooth proper family of geometrically connected curves of fixed genus over a reduced base.
\end{lemma}
\begin{proof}
For a semilinear endomorphism $F$ of a $g$-dimensional vector space over a perfect field, the subspaces $F^j(V)$ decrease. Once two consecutive subspaces agree, all subsequent ones agree. There can be at most $g$ strict dimension drops, so the stable image is $F^g(V)$. The asserted locus is therefore defined by nonvanishing of minors of the displayed bundle map.

For the curve assertion, put $\mathcal H=R^1f_*\OO_{\mathcal C}$. Its local freeness and base change follow from perfect cohomology and the constant fibre dimensions $h^0=1$, $h^1=g$, $h^i=0$ for $i>1$. Here is the elementary passage from these facts to local freeness. Represent the perfect complex locally by a bounded complex of finite free modules. Starting at its lowest term, the constant cohomology dimensions determine the constant ranks of the differentials successively. Over a reduced base, a matrix of constant fibre rank has locally an invertible maximal minor and zero complementary block: its larger minors vanish at every point and therefore vanish as functions. Kernels, images and cokernels are locally free, and the complex splits into its cohomology and contractible summands locally. This also proves base change. Relative Frobenius gives the indicated map on $\mathcal H$, and stable coherent Frobenius rank is the curve's $2$-rank.
\end{proof}

\begin{proof}[Proof of Theorem~\ref{thm:eventual-all-mixed}]
Choose a finitely generated field $K\subset k$ over $\FF_2$ over which $(A,M)$ is defined. After spreading and shrinking, there is an integral finite-type scheme $S/\FF_2$ with function field $K$, an abelian scheme $\pi:\mathcal A\to S$ of relative dimension two, and a relatively ample line bundle $\mathcal M$. Its geometric generic fibre, extended to $k$, is the original pair. The integer $\chi(\mathcal M_s)$ is the constant $D$.

We first arrange that every geometric fibre is mixed. By perfect cohomology, followed by splitting the complex at the generic point and shrinking, $\mathcal E=R^1\pi_*\OO_{\mathcal A}$ is a rank-two vector bundle commuting with base change. Frobenius gives
\[
 \Phi:F_S^*\mathcal E\longrightarrow\mathcal E.
\]
Its second iterate has rank one at the geometric generic point, by the preceding lemma and $f_A=1$. All its two-by-two minors vanish on integral $S$. Some entry is nonzero generically; restrict to a nonempty trivializing open where that entry is invertible. The second iterate now has rank one at every geometric point. Every fibre therefore has $2$-rank one.

Choose a closed point $s$ of this fixed nonempty finite-type scheme. Its residue field is finite, and the actual pair $(A_s,M_s)$ is defined over that field. Let $a_s=a(M_s)\ge0$. If $a_s>0$, Theorem~\ref{thm:finite-field} supplies an integer $e_s$ such that, for every $e\ge e_s$, the geometric system $|M_s^{3e}|$ has a smooth curve of defect two. If $a_s=0$, Theorem~\ref{thm:p10exact}(1) supplies a smooth curve of defect at most two for every $e\ge1$; take $e_s=1$. In both cases, for every $e\ge e_s$ there is a smooth curve $C_s$ satisfying
\begin{equation}\label{eq:special-rank-bound}
                    f_{C_s}\ge 9De^2-1.
\end{equation}
This argument does not require constancy of the ordinary index under specialization.

Relative Serre vanishing on the fixed Noetherian base gives an integer $e_1$ such that, for all $e\ge e_1$, the higher direct images of $\mathcal M^{3e}$ vanish and
\[
                  \mathcal E_e=\pi_*\mathcal M^{3e}
\]
is locally free and commutes with arbitrary base change. Enlarge $e_1$ so that these powers are relatively generated. This is one threshold on the fixed quasi-compact base, and it does not require deleting the chosen point $s$.

Let $P_e$ be the projective bundle of one-dimensional spaces of sections of $\mathcal E_e$. It is integral, and each geometric fibre is the full complete linear system of $M_s^{3e}$. Its tautological section defines a universal divisor. On each smooth integral abelian fibre a nonzero section is a nonzerodivisor; hence this is a relative effective Cartier divisor, flat over $P_e$ \citep[Tag~062Y]{Stacks}. Let $V_e\subset P_e$ be its smooth-fibre locus. This locus is open because the nonsmooth locus in the proper universal divisor has closed image. Its fibres are geometrically connected: on a geometric abelian fibre $A_t$ over an algebraically closed field $\Omega$, put $N=\mathcal M_t^{3e}$. Lemma~\ref{lem:cohomology}, or the sequence
\[
 0\longrightarrow N^{-1}\longrightarrow\OO_{A_t}
       \longrightarrow\OO_C\longrightarrow0,
\]
together with $H^0(N^{-1})=H^1(N^{-1})=0$, gives $H^0(\OO_C)=\Omega$. The genus is $g_e=9De^2+1$ by adjunction and Riemann--Roch.

Fix any $e\ge\max(e_s,e_1)$. The curve in~\eqref{eq:special-rank-bound} gives a geometric point of $V_e$ where $2$-rank is at least $g_e-2$. Lemma~\ref{lem:stable-rank-open} shows that this inequality defines a nonempty open $W_e\subset V_e$. Because $P_e$ is integral, $W_e$ contains its generic point and has a nonempty generic fibre over $S$. Extending that geometric generic fibre to $k$ gives a nonempty open in the original complete system $|M^{3e}|$ with
\[
                             f_C\ge g_e-2.
\]
The positive-index upper bound in Theorem~\ref{thm:p10bounds} gives the reverse inequality, since $a(M)>0$. Thus $f_C=g_e-2$ generically, and~\eqref{eq:dimensions} gives $r(C)=9De^2-2$.

The same point $s$ and the same two integers $e_s,e_1$ serve every $e\ge e_0=\max(e_s,e_1)$. Only the projective bundle and its open subset depend on $e$. There is no intersection of infinitely many open subsets of $S$. Since the original pair is the geometric generic fibre of its own spread, the conclusion holds for every pair in the theorem.
\end{proof}

The proof is existential for an arbitrary pair. For a pair already defined over a finite field, Theorem~\ref{thm:finite-field} retains its explicit sufficient bound. An index-zero ample bundle on the special fibre need not extend to the generic fibre; the proof transfers a nonvanishing Frobenius minor instead. Thus eventual sharp rank does not imply existence of an index-zero ample class on the original surface.

\section{A mixed surface with no index-zero ample class}
\label{sec:no-zero-class}

This section proves that the class-existence hypothesis in the ample-difference theorem fails on some mixed surfaces, and computes every ample polarization index on an explicit isogeny construction.

The construction starts from a geometric generic Jacobian with full symplectic monodromy. The passage from this monodromy to endomorphism ring $\mathbf Z$ is classical; compare \citet[Corollary~4.2, p.~297]{Chai}. Quotients of Picard-rank-one principally polarized varieties already supply classical obstructions to principal polarizations \citep[\S11.24, p.~169]{EGM}. We give the integral ordinary-index calculation, which is the feature needed here.

\begin{lemma}\label{lem:mixed-principal-rank-one}
There is an algebraically closed field $k$ of characteristic two and a principally polarized mixed abelian surface $(B,\Theta)/k$ such that
\[
             \End(B)=\mathbf Z,\qquad \NS(B)=\mathbf Z[\Theta].
\]
\end{lemma}
\begin{proof}
The smooth projective Artin--Schreier curve with function field
\[
                         y^2+y=x^3+\frac1x
\]
has genus two and $2$-rank one. Indeed its two poles have odd orders three and one. Their different contributions are four and two, so Riemann--Hurwitz gives $2g-2=-4+4+2=2$; the two-pole $2$-rank formula gives $f=2-1=1$. These formulas, with their hypotheses and proofs, are \citet[Lemmas~2.6--2.7 and Example~2.8, pp.~711--712]{PriesZhu}. In particular the $2$-rank-one stratum of the genus-two moduli space is nonempty.

Choose an irreducible component of this stratum over $\overline{\FF}_2$. Pass to a connected component with full level-five structure and then to a connected normal open scheme carrying the universal curve and its principally polarized Jacobian. The geometric monodromy theorem of \citet[Theorem~4.5, pp.~14--15 of the author version; scheme rigidifications in \S5.1, pp.~17--18]{AchterPries} applies in characteristic two, genus two and $2$-rank one. Before the finite level cover its $3$-adic image is $\operatorname{Sp}_4(\mathbf Z_3)$. A finite cover changes this to a finite-index subgroup; restriction to a nonempty open does not decrease the image, since a connected finite \'etale cover of a normal integral scheme remains connected over that open. Its Zariski closure is therefore $\operatorname{Sp}_{4,\mathbf Q_3}$. Let $B$ be the Jacobian at the geometric generic point, over an algebraic closure $k$ of the function field. Its canonical principal polarization is represented by an ample line bundle $\Theta$, and $B$ is mixed.

We spell out why the geometric endomorphism ring has rank one. It is a finite free abelian group, and all its generators are defined over a common finite extension of the generic field \citep[Theorem~12.10 and Corollaries~12.11--12.13, pp.~182--184]{EGM}. For the finite-generation step, the classical positive integral form on the endomorphism group gives an integer Gram matrix on a rational basis. Its nonzero determinant bounds all denominators of every endomorphism, so the group lies in a fixed finite-rank lattice. After spreading out and shrinking, the generators act on the abelian scheme over the normalization in that extension. A purely inseparable extension leaves the \'etale fundamental group unchanged, and the separable part becomes finite \'etale after shrinking. Thus every generator commutes with a finite-index monodromy subgroup, still Zariski dense in $\operatorname{Sp}_4$.

The centralizer of $\operatorname{Sp}_4$ on its defining representation is $\mathbf Q_3$ times the identity. Commutation with the diagonal torus makes an endomorphism diagonal in the four standard weight spaces; the elementary symplectic root matrices force its four diagonal entries to agree. The Tate-module map is injective after tensoring $\End(B)$ with $\mathbf Z_3$: an endomorphism killing $B[3^m]$ factors through the quotient $[3^m]:B\to B$, so its kernel modulo $3^m$ is exactly $3^m\End(B)$. Taking inverse limits, using finite generation, gives the injection; this is also \citet[Lemmas~12.8--12.9, pp.~182--183]{EGM}. Hence
\[
      \End(B)\otimes\mathbf Q_3\ \subset\ \mathbf Q_3\operatorname{id}.
\]
It follows that $\End^0(B)=\mathbf Q$. A finite $\mathbf Z$-algebra contained in $\mathbf Q$ consists of rational algebraic integers, so $\End(B)=\mathbf Z$.

Finally $L\mapsto\lambda_\Theta^{-1}\lambda_L$ injects $\NS(B)$ into $\End(B)$: the kernel on $\Pic(B)$ is exactly $\Pic^0(B)$ \citep[\S\S7.22--7.26, pp.~106--108]{EGM}. It sends $[\Theta]$ to the identity. Therefore $\NS(B)=\mathbf Z[\Theta]$.
\end{proof}

\begin{proof}[Proof of Theorem~\ref{thm:no-ordinary-index-zero-ample-class}]
Choose $(B,\Theta)$ as in the lemma. The local-local height-two summand of $B[2]$ contains a subgroup $H\simeq\alpha_2$. More explicitly, the mixed surface has
\[
 B[2]\simeq(\mathbf Z/2\mathbf Z\oplus\mu_2)\oplus I_{1,1},
\]
where $I_{1,1}$ is the $2$-torsion of a supersingular elliptic curve and contains a unique $\alpha_2$; see \citet[Example~2.2, Lemma~3.1 and Propositions~3.2--3.3, author version pp.~1--4; genus-two table p.~6]{PriesTorsion}. Naturality of relative Frobenius gives $H\subset\ker(F_B:B\to B^{(2)})$.

Form the quotient isogeny $q:B\to A=B/H$. Its degree is two and it is radicial. The relative Frobenius, of degree four, factors as
\[
                    F_B=hq,\qquad h:A\longrightarrow B^{(2)},
\]
where $h$ is an isogeny of degree two. Define the actual line bundle
\begin{equation}\label{eq:no-zero-pullback}
 L=h^*(\Theta^{(2)}),\qquad q^*L\simeq F_B^*(\Theta^{(2)})
                              \simeq\Theta^{\otimes2}.
\end{equation}
It is ample because $h$ is finite. The isogeny preserves the $2$-rank, so $A$ is mixed. Riemann--Roch and pullback by the degree-two isogeny give
\[
                       2\chi(L)=\chi(2\Theta)=4,
                       \qquad \chi(L)=2.
\]

The kernels of both $q$ and $q^\vee$ are $\alpha_2$, since this group is Cartier self-dual. Their maps on integral geometric $2$-adic Tate modules are isomorphisms. Indeed a radicial isogeny is bijective on algebraically closed points; if its image of $x$ is killed by $2^m$, then $2^mx$ is a point of its infinitesimal kernel and must be zero. Thus the bijection restricts to every geometric $2^m$-torsion group. The polarization identity from \eqref{eq:no-zero-pullback} is
\[
                         q^\vee\lambda_Lq=2\lambda_\Theta.
\]
The outer Tate-module maps are integral isomorphisms and $\lambda_\Theta$ is principal. On the rank-one Tate modules the middle map therefore has cokernel length one. Consequently
\[
                         a(L)=1,\qquad b(L)=v_2(2)-1=0.
\]

Pullback by $q$ injects N\'eron--Severi groups: its kernel is killed by the degree through the norm identity, whereas $\NS(A)$ is torsion-free by its injection into $\Hom(A,A^\vee)$. For any line bundle $N$ on $A$, write $[q^*N]=m[\Theta]$. Euler characteristic depends only on this class, and
\[
                            2\chi(N)=m^2.
\]
Thus $m$ is even. Equation~\eqref{eq:no-zero-pullback} and injectivity imply $[N]=(m/2)[L]$. This proves $\NS(A)=\mathbf Z[L]$. If $N$ is ample then $t=m/2>0$, because $q^*N$ is ample. Algebraically trivial twists do not change the polarization homomorphism, so
\[
 \lambda_N=t\lambda_L,\qquad \chi(N)=2t^2,
 \qquad a(N)=1+v_2(t),\qquad b(N)=v_2(t).
\]
In particular no ample $N$ has ordinary index zero. Theorem~\ref{thm:finite-field} also shows that this $A$ cannot have a finite-field model.
\end{proof}

The same descent follows from the following theta-group calculation, an elementary special case of the classical finite-group extension statement \citep[Fact~8.7, p.~115]{EGM}.

\begin{lemma}[$\alpha_2$ theta descent]\label{lem:alpha-two-theta-descent}
Over an algebraically closed field of characteristic two, every central extension of fppf group schemes
\[
                  1\longrightarrow\mathbf G_m\longrightarrow E
                   \longrightarrow\alpha_2\longrightarrow1
\]
splits. If $H\simeq\alpha_2$ is a subgroup of the polarization kernel $K(P)$ of a line bundle $P$ on an abelian variety, then $P$ descends to the quotient by $H$.
\end{lemma}
\begin{proof}
The underlying $\mathbf G_m$-torsor is trivial because $\alpha_2=\operatorname{Spec}k[t]/(t^2)$ has local coordinate ring and hence trivial Picard group. Choose a section taking zero to the identity. Centrality makes multiplication on the resulting torsor take the form
\[
        (u,x)(v,y)=(uvc(x,y),x+y),\qquad
 c\in\bigl(k[x,y]/(x^2,y^2)\bigr)^\times.
\]
The identity conditions give $c(x,0)=c(0,y)=1$, so $c(x,y)=1+\gamma xy$ for some $\gamma\in k$. In particular it is symmetric. Choose $\beta^2=\gamma$ and set $d(t)=1+\beta t$. Direct multiplication in this coordinate ring gives
\[
                \frac{d(x)d(y)}{d(x+y)}=1+\beta^2xy=c(x,y),
\]
since $(x+y)^2=x^2=y^2=0$. Multiplying the chosen section by $d(t)^{-1}$ makes its cocycle one. It is then a group homomorphism and splits the extension.

Restrict the theta-group extension of $P$ to $H$. Its splitting is precisely an $H$-linearization of $P$. Faithfully flat descent along the quotient torsor gives the required line bundle; see \citet[Theorem~8.10 and its proof, p.~116]{EGM}. In the construction above $H\subset B[2]=K(2\Theta)$, so this also descends $\Theta^{\otimes2}$ and yields the same polarization-index calculation.
\end{proof}


\section{The toral image and a double-cover deformation}
\label{sec:toral}

We now give an exact reduction of the missing sharpness assertion. Its hypotheses include an already sharp source. The reduction is a sufficient inductive mechanism, not a proof that every sharp curve must arise from this mechanism.

Fix a mixed surface $B$ and an ample bundle $H=M^{3e}$ whose ordinary index is positive. Write $m=\chi(H)=h^0(H)$. The multiplicative subgroup $\mu _2\subset B[2]$ is isotropic for $H$ and lies in $\ker\lambda_H$; the descent theorem used above linearizes $H$ under this subgroup. Its infinitesimal generator acts on $W=H^0(B,H)$ as an idempotent $e$, with weight decomposition $W=W_0\oplus W_1$. Changing a linearization exchanges the two weights and changes none of the constructions below. Fix a nonzero invariant volume form on $B$.

For $s\in W$, coherent Frobenius is the semilinear map
\[
 F_s=sF:H^2(B,H^{-1})\longrightarrow H^2(B,H^{-1}).
\]
Here $F$ first lands in $H^2(B,H^{-2})$, and multiplication by $s$ returns to $H^2(B,H^{-1})$. Serre duality identifies its dual, with the relative Frobenius twists understood, with Cartier multiplication by $s$. We write $\beta_F(s,t)=C(st\omega)$ for this Cartier pairing. It is perfect on the Frobenius pushforward, alternating because $C(\omega)=0$, and orthogonal for the two weights. Indeed $C(e(f)\omega)=0$ and the Leibniz rule give
\[
 \beta_F(es,t)=\beta_F(s,et),\qquad
 \beta_F(s,s)=\beta_F(s,es)=0.
\]
The local Cartier formula in two regular parameters proves these assertions directly; its duality convention is also used in Lemma~\ref{lem:tensor}.

Define
\[
 I_s=\operatorname{ann}\langle s,es\rangle\subset H^2(B,H^{-1}).
\]
Call $s$ \emph{sharp} if $F_s$ is bijective on $I_s$. Such a section has relative stable rank $m-2$ and curve defect two. The condition is open.

\begin{lemma}\label{lem:toral-image}
One has $\operatorname{im}F_s\subset I_s$. The global kernel sheaf $\mathcal I$ of the two evaluation functionals is reflexive, has $H^0(\mathcal I)=0$, and has determinant $\OO(-m)$. Away from
$\mathbf P(W_0)\cup\mathbf P(W_1)$ it is the rank-$m-2$ bundle with fibres $I_s$. Frobenius on this bundle has determinant a section of $\OO(m)$, extending across the excluded locus. Its nonvanishing is equivalent to sharpness on that open. Every sharp section lies there.
\end{lemma}
\begin{proof}
The two displayed Cartier identities are dual to the image containment. On $\mathbf P(W)$ consider the map
\[
 W^*\otimes\OO(-1)\longrightarrow\OO^2,\qquad
 \ell\longmapsto(\ell(s),\ell(es)).
\]
Its rank is two precisely off the two weight spaces. Its image sheaf is torsion-free; the depth lemma therefore makes its kernel reflexive on the smooth parameter space. Its injection into $W^*\otimes\OO(-1)$ gives $H^0(\mathcal I)=0$. Each weight space has codimension $m/2\ge2$, as follows from Lemma~\ref{lem:two-bundles}. The determinant of the kernel on the complement is $\OO(-m)$, and this extends uniquely across codimension at least two. Hence the determinant of $F^*\mathcal I\to\mathcal I$ belongs to
$\det\mathcal I\otimes(\det F^*\mathcal I)^{-1}=\OO(m)$ and extends there as a section of a line bundle on the normal parameter space.
The full map $F_s$ has rank at most $m-2$ on that dense open, and therefore everywhere, since its $(m-1)$-minors vanish identically. At a weight section $\dim I_s\ge m-1$, so its restriction cannot be bijective. The quotient of the full cohomology space by $I_s$ has zero Frobenius; thus the stable rank of the full map is the stable rank on $I_s$. The cohomology sequence for a divisor identifies its unit part with that of the relative quotient because Frobenius on $H^2(B,\OO_B)$ is zero. This proves the assertions about curve rank.
\end{proof}

\begin{lemma}\label{lem:double}
Let $\pi:\widetilde B\to B$ be a connected \'etale double isogeny and let $s$ be sharp. Put $\widetilde H=\pi^*H$. On $I_{\pi^*s}$ the first and stable Frobenius ranks are both $2m-4$; its kernel and cokernel have dimension two. The linear first variation from kernel to cokernel vanishes in all pullback directions. It therefore factors through the section trace $H^0(\widetilde H)\to H^0(H)$. It also kills $\langle s,es\rangle$ in this trace space.
\end{lemma}
\begin{proof}
The Artin--Schreier trace sequence is
\[
 0\to\OO_B\to\pi_*\OO_{\widetilde B}\xrightarrow{\operatorname{Tr}}\OO_B\to0.
\]
After tensoring by $H^{-1}$, ample vanishing gives a Frobenius-compatible exact sequence
\[
 0\to V\longrightarrow\widetilde V\xrightarrow{\operatorname{Tr}}V\to0,
 \qquad V=H^2(H^{-1}),\quad\widetilde V=H^2(\widetilde H^{-1}).
\]
Compatibility follows locally from $\operatorname{Tr}(x^2)=\operatorname{Tr}(x)^2$. Trace is dual to section pullback, so $I_{\pi^*s}=\operatorname{Tr}^{-1}(I_s)$. Thus
\[
 0\to V\longrightarrow I_{\pi^*s}\xrightarrow{\operatorname{Tr}}I_s\to0.
\]
Since $\operatorname{im}F_s\subset I_s$ and $F_s|I_s$ is bijective, its full image equals $I_s$ and its kernel $K$ has dimension two, disjoint from $I_s$. The kernel--cokernel sequence, applied after Frobenius linearization, identifies the kernel on $I_{\pi^*s}$ with $\pi^*K$ and its cokernel with $V/I_s$. If an image element equals $\pi^*k$ with $k\in K$, taking trace forces its preimage to have trace zero, hence to come from $V$. Then $k\in K\cap I_s=0$. The first image is therefore already bijective under Frobenius, and its dimension is $2m-4$.

The tangent map of a family of linear maps of constant first rank is zero from kernel to cokernel. The sharp source open gives such a family inside the pullback section subspace, by the preceding argument. Thus invariant directions have zero variation. The trace sequence for $H$, together with $H^1(H)=0$, identifies the quotient by pullback directions with $H^0(H)$.

For the last assertion use compatible Cartier-dual spaces. Write $s=s_0+s_1$ by weight. Sharpness says
$W=\langle s_0,s_1\rangle\oplus\operatorname{im}C_s$.
Choose weight-preserving trace lifts $q_i$ of $s_i$. Correct them by pullback sections $w_i$ until $C_{\pi^*s}(q_i+\pi^*w_i)$ lies in the pullback of $\langle s_0,s_1\rangle$. Their classes give the dual cokernel basis of the restricted map. For $k\in K$, the variation coefficient in a direction $t$ is
\[
 \big\langle k,C_B(\operatorname{Tr}(t(q_j+\pi^*w_j))\omega)\big\rangle.
\]
For $t=q_i$, the correction term is $\beta_F(s_i,w_j)$ and belongs to $\operatorname{im}C_s$. The other term is zero: for different weights the pairing is zero, and for equal weights it is the Cartier image of a square times $\omega$. Thus both $s_i$ are killed in trace coordinates.
\end{proof}

An invertible member of this two-by-two variation family proves sharpness on a nonempty open of the cover's full section space. To see sufficiency, split the linearized matrix at the pullback point into an invertible $(2m-4)$-block and its two-dimensional kernel and cokernel. Its determinant has lowest potentially nonzero term equal, up to a nonzero scalar, to the determinant of that first variation. No converse about arbitrary higher-order deformations is asserted.

\section{Two bundles and a determinant curve}
\label{sec:curve}

This section constructs the two Cartier bundles and their determinant curve, including the torsor and pencil exact sequences.

Factor relative Frobenius as
\[
 B\xrightarrow{\rho}C=B/\mu _2\xrightarrow{\nu}T=B^{(2)}.
\]
The map $\nu$ is an $\alpha _2$-torsor. Write $H=\rho^*J$, where $J$ is a cube of an ample bundle. The torsor algebra of $\rho$ has the form $\OO_C\oplus Q_0$, where $Q_0$ is the nontrivial geometric two-torsion line bundle. The dual kernel of $\nu$ is $\alpha _2$, so $\nu^*$ is bijective on geometric points of $\Pic^0$. Hence $Q_0=\nu^*P$ for the unique nontrivial $P\in\Pic^0(T)[2](k)$. Fix its induced trivialization $P^2\simeq\OO_T$. Put
\[
 L=H^{(2)},\quad V=\nu_*J,\quad E_0=V,\quad E_1=VP,
 \qquad \mathcal L=LP.
\]

\begin{lemma}\label{lem:two-bundles}
The decomposition $F_{B/k,*}H=E_0\oplus E_1$ is orthogonal for the perfect alternating Cartier pairing with values in $L$. In particular $\det E_i=L$. The bundles $E_i$ are globally generated,
\[
 h^0(E_i)=c_2(E_i)=m/2,\qquad H^{>0}(E_i)=0,
 \qquad H^\bullet(P)=H^\bullet(\operatorname{Hom}(E_0,E_1))=0.
\]
\end{lemma}
\begin{proof}
Projection formula gives the displayed decomposition, and the Cartier identities already proved give orthogonality and perfectness on each rank-two summand. A perfect alternating pairing in rank two identifies its determinant with the pairing line. Ample vanishing on $C$ gives higher cohomology zero and $h^0(E_i)=\chi(J)=m/2$. Riemann--Roch on the abelian surface gives
$\chi(E_i)=c_1(L)^2/2-c_2(E_i)=m-c_2(E_i)$.

The vanishing of every cohomology group of a nontrivial algebraically trivial line bundle is \citet[Lemma~7.19 and proof, p.~105]{EGM}. For completeness, if $H^0(P)\ne0$, the same holds for $P^{-1}=(-1)^*P$, so $P$ is trivial. If $i>0$ is the least degree with nonzero cohomology, the factorization of the identity through addition $T\to T\times T\to T$ and $m^*P=P\boxtimes P$ contradicts K\"unneth: every summand in degree $i$ has a smaller nonzero degree. This proves acyclicity.

The diagonal ideal in $C\times_T C$ is square-zero and trivial as a line on $C$. Therefore $\nu^*\nu_*J$ has a two-step filtration with quotients $J$. Finite-flat base change and projection formula identify
$H^\bullet(E_0^\vee E_1)$ with the cohomology of
$(\nu^*\nu_*J)^\vee JQ_0$, whose two quotients are $Q_0$. They are acyclic by the preceding argument.

Every cube of an ample bundle on an abelian variety is very ample \citep[Theorem~2.2.3, p.~13]{GenestierNgo}. The same holds after a degree-zero twist, by taking a cube root of that twist. Hence $J$ and $JQ_0$ separate the length-two fibres of $\nu$, proving global generation of their pushforwards.
\end{proof}

Choose a general section pair $(s_0,s_1)\in H^0(E_0)\oplus H^0(E_1)$. Its zero schemes $Z_0,Z_1$ are finite and disjoint: global generation makes simultaneous vanishing impose four conditions over a surface, and makes a general section of each rank-two bundle have zero locus of dimension at most zero. Each zero scheme has length $m/2$. The contraction
\[
 E_0\oplus E_1\longrightarrow L,\qquad
 (t_0,t_1)\longmapsto\beta_F(s_0,t_0)+\beta_F(s_1,t_1)
\]
is surjective. Write $K_s$ for its kernel. The map
$P\to K_s$, $r\mapsto(rs_1,rs_0)$ is a subbundle; put $Q_s=K_s/P$.
The two vectors $(s_0,0)$ and $(0,s_1)$ descend to a map
$a:\OO_T^2\to Q_s$.

\begin{lemma}\label{lem:curve-sequences}
Put $w=s_0\wedge s_1\in H^0(T,\mathcal L)$ and
$D_s=(w=0)$. This is a nonzero Cartier divisor of arithmetic genus $m+1$, and
\[
 \det Q_s=\mathcal L,\qquad
 0\to\OO_T^2\xrightarrow{a}Q_s\to i_*R_s\to0.                 \tag{C1}
\]
The cokernel is a line bundle on $D_s$, even when this divisor is nonreduced. The common image of the two sections on $D_s$ is a line subbundle $R_0\subset V|_{D_s}$. With $A_s=R_0^2$, one has
\[
 R_0=\OO_{D_s}(Z_0),\quad R_0P=\OO_{D_s}(Z_1),\quad
 R_s=\omega_{D_s}A_s^{-1},\quad\deg A_s=m.                  \tag{C2}
\]
The squared sections $a_0,a_1$ form a base-point-free pencil in $A_s$. If $s$ is sharp, then
\[
 h^0(Q_s)=h^1(Q_s)=h^0(A_s)=h^1(A_s)=h^0(R_s)=2,\qquad h^2(Q_s)=0.
                                                                    \tag{C3}
\]
\end{lemma}
\begin{proof}
The contraction kernel has determinant $L$, and quotienting by $P$ gives $\det Q_s=LP$. In a frame with $s_0=(1,0)$ write $s_1=(h,w)p$, where $p$ frames $P$ and $p^2=u$ is an arbitrary unit. Solving the contraction equation and quotienting by the cross vector gives
\[
                       a=\begin{pmatrix}1&uh^2\\0&w\end{pmatrix}.
                                                                    \tag{C4}
\]
Thus the determinant is $w$, and the cokernel along $D_s$ is locally free of rank one. A general pair has $w\ne0$, since global generation permits independent values at a point. The class $\mathcal L$ is ample and $\chi(\mathcal L)=m$, proving its genus assertion.

On $D_s$ at least one section is nonvanishing, so the rank-one image is a subbundle. Near $Z_0$ its vanishing ideal on the surface is $(h,w)$, a regular sequence; hence $h$ is a nonzerodivisor on $D_s$. This proves the Cartier-divisor and line-bundle assertions for $R_0$, and likewise for $Z_1$. Their squares have no common zero. Tensoring (C1) with $\OO_{D_s}$ gives the Tor kernel $R_s(-D_s)$. In (C4) this is the kernel of the pencil row $(1,uh^2)$, namely $A_s^{-1}$. Adjunction gives $R_s=\omega_{D_s}A_s^{-1}$.

Sharpness gives two-dimensional kernel and cokernel for global Cartier contraction. The acyclicity of $P$ identifies these with $H^0(Q_s)$ and $H^1(Q_s)$; higher cohomology vanishes. The assertions for the pencil will follow from the next lemma, whose proof does not use them.
\end{proof}

\begin{lemma}[The norm pencil and both torsor classes]\label{lem:norm-pencil}
There is a nonzerodivisor $\sigma\in H^0(D_s,R_sP)$ such that the norm sections $S_i\in H^0(T,L)$ satisfy
\[
 S_i|_{D_s}=\sigma a_i,\qquad
 H^0(T,L)/\langle S_0,S_1\rangle
  =H^0(D_s,\omega_{D_s}P)/\sigma\langle a_0,a_1\rangle.       \tag{C5}
\]
There is an exact curve sequence
\[
 0\to A_s\to V|_{D_s}\otimes R_0\to L|_{D_s}\to0,            \tag{C6}
\]
whose global map is the entire Cartier contraction and whose middle term has $H^1=0$. If $\zeta\in H^1(T,\OO_T)$ is the $\alpha _2$-torsor class of $\nu$, its boundary sends $\sigma a$ to $\zeta|_{D_s}\cup a$. At a sharp section this cup map on $H^0(A_s)\to H^1(A_s)$ is an isomorphism.

Let $\alpha$ be the coherent class of a connected \'etale double cover of $T$. The full Cartier contraction after that pullback has kernel dimension
\[
       2+\dim\ker[\alpha\cup-:H^0(A_s)\to H^1(A_s)].         \tag{C7}
\]
This full first-rank statement is distinct from Lemma~\ref{lem:double}.
\end{lemma}
\begin{proof}
Locally write the torsor algebra as $\OO_T[z]/(z^2-h)$. The norm of $a+bz$ is $a^2+hb^2$, with target $\det V=L$, since $\det\nu_*\OO_C$ is trivial. On $VP$ the norm is multiplied by the local unit $u$ defining $P^2$. Its restriction to the line $R_0$ is a map $A_s\to L|_{D_s}$, giving $\sigma$ and (C5). For a general pair the two divisors of $J,JQ_0$ have no common component. Since $\nu$ is a homeomorphism, their norm divisors do not either. If $\sigma$ vanished at a minimal prime of $D_s$, both norms would contain that component. A Cartier divisor on a smooth surface is Cohen--Macaulay, with no embedded associated primes, so $\sigma$ is a nonzerodivisor. Restriction $H^0(L)\to H^0(L|_{D_s})$ is an isomorphism by the acyclic kernel $P$, proving the quotient formula.

The original pair has the surface presentation
\[
 0\to\OO_T\oplus P\to V\to i_*(L|_{D_s}R_0^{-1})\to0.
\]
Duality for a Cartier divisor gives
$0\to V^\vee\to\OO_T\oplus P\to i_*(R_0P)\to0$.
After tensoring by $VP$, the left term is acyclic by Lemma~\ref{lem:two-bundles}, and the middle term is $E_1\oplus E_0$. Thus its right term $V|_{D_s}R_0$ has the same global sections and no higher first cohomology. Wedge with the common line gives (C6), and its composite from $E_0\oplus E_1$ is exactly the Cartier contraction. Consequently its kernel and cokernel are $H^0(A_s),H^1(A_s)$. At a sharp section they have dimension two; curve duality gives the remaining part of (C3).

For the boundary formula choose torsor coordinates $z_j-z_i=c_{ij}$. If $r$ frames $R_0$, then $r\wedge z_ir=n(r)=\sigma r^2$. A local lift of $\sigma(fr^2)$ in (C6) is $f(z_ir)\otimes r$, and its difference on overlaps is $c_{ij}fr^2$. This proves the cup formula. Sharpness is exactly that the norm-section complement maps isomorphically to the Cartier cokernel, hence that this cup map is invertible.

All sequences commute with \'etale base change. The required acyclicity persists, since pushforward along a double cover is an extension of two copies. Therefore the pulled-back kernel is $H^0(\widetilde D_s,A_s|_{\widetilde D_s})$. The trace sequence on $D_s$ has boundary $\alpha\cup-$, giving
$0\to H^0(A_s)\to H^0(A_s|_{\widetilde D_s})\to\ker(\alpha\cup-)\to0$.
This proves (C7).
\end{proof}

The line bundle $P$ also carries an invariant logarithmic differential. In a frame with $p^2=u$, the forms $d\log u$ glue to $\eta_P$ and satisfy $C(\eta_P)=\eta_P$. They are not all zero: otherwise every $u$ would locally be a square on the smooth perfect-base surface, and its unique square roots would trivialize $P$. Formula (C4) gives
\[
                       d\log(a_1/a_0)=\eta_P|_{D_s}.       \tag{C8}
\]
This statement does not yet say that the differential on the possibly nonreduced curve is nonzero.

\section{An alternating factorization and trace-zero variation}
\label{sec:tracezero}

This section proves trace zero through an acyclic alternating factorization and identifies the Frobenius-typed first variation.

Set $\mathfrak g=(\operatorname{End}^0(V)/\OO_T)P$, where
$\operatorname{End}^0(V)$ is the trace kernel. Scalar matrices lie in this kernel in characteristic two.

\begin{lemma}\label{lem:factor}
The bundle $\mathfrak g$ is acyclic and has a perfect alternating trace pairing. Let $J_{\mathfrak g}:\mathfrak g^\vee\to\mathfrak g$ be its inverse and $J_2=\left(\begin{smallmatrix}0&1\\1&0\end{smallmatrix}\right)$. There is a map $f:\mathfrak g\to Q_s$ such that
\[
 aJ_2a^\vee=fJ_{\mathfrak g}f^\vee,\qquad
 0\to Q_s^\vee\xrightarrow{(a^\vee,f^\vee)}
 \OO_T^2\oplus\mathfrak g^\vee
 \xrightarrow{(aJ_2,fJ_{\mathfrak g})}Q_s\to0.              \tag{F1}
\]
Its first image is Lagrangian. At a sharp section it induces
\[
 0\to H^1(Q_s^\vee)\to H^1(\OO_T)^2\to H^1(Q_s)\to0.       \tag{F2}
\]
Moreover $\operatorname{coker}f=i_*A_s$.
\end{lemma}
\begin{proof}
Both $P$ and $\operatorname{End}(V)P$ are acyclic; the latter is $\operatorname{Hom}(E_0,E_1)$. Apply cohomology to
\[\begin{split}
0&\to\operatorname{End}^0(V)P\to\operatorname{End}(V)P\to P\to0,\\
0&\to P\to\operatorname{End}^0(V)P\to\mathfrak g\to0.
\end{split}\]
This proves acyclicity. Modulo scalars, a trace-zero matrix is represented by its off-diagonal entries $(b,c)$; the trace pairing is $bc'+cb'$, alternating and perfect.

For $\phi\in\operatorname{End}^0(V)P$, the pair
$(\phi s_1,\phi s_0)$ belongs to $K_s$, because its contraction is
$\operatorname{tr}(\phi)(s_0\wedge s_1)=0$. Scalars give its existing subline $P$, so the map descends to $f$. In the chart (C4),
\[
 a=\begin{pmatrix}1&uh^2\\0&w\end{pmatrix},\quad
 f=\begin{pmatrix}uw&0\\0&1\end{pmatrix},\quad
 J_{\mathfrak g}=u^{-1}J_2.
\]
Both sides of the first identity in (F1) are $wJ_2$. These are identities without choosing a square root of $u$. The images of $a$ and $f$ generate $Q_s$, so the last arrow in (F1) is surjective. The zero composite, ranks and perfect alternating form on the middle bundle identify its kernel with the first image and prove the Lagrangian assertion. The local cokernel of $f$ is the first-coordinate line on $D_s$, with its two sections $(1,uh^2)$; hence it is $A_s$.

At a sharp section the map $H^0(\OO^2\oplus\mathfrak g^\vee)\to H^0(Q_s)$ is an isomorphism. Its Serre dual is the corresponding isomorphism in degree two. The remaining long cohomology sequence is exactly (F2).
\end{proof}

\begin{theorem}[Trace zero and the remaining singular line]\label{thm:tracezero}
At a sharp section the torsor cup matrix
\[
 T_s=(\zeta\cup-)^{-1}(\alpha\cup-)
       \in\operatorname{End}H^0(D_s,A_s)
\]
has trace zero. Every actual trace-direction first variation in Lemma~\ref{lem:double} has trace zero under the compatible kernel and cokernel identifications. If all these variations are singular, their image is zero or a single line generated by a rank-one nilpotent matrix. An image of dimension at least two is therefore sufficient for a sharp deformation; its existence is not asserted.
\end{theorem}
\begin{proof}
Dualizing (C1) gives
\[
 0\to Q_s^\vee\to\OO_T^2\to i_*A_s\to0.
\]
At sharpness its last row is the basis $(a_0,a_1)$ of global sections. Thus the image $S$ of $H^1(Q_s^\vee)$ in $H^1(\OO_T)^2$ is the kernel of the cup map
$(\gamma_0,\gamma_1)\mapsto\gamma_0a_0+\gamma_1a_1$.
The cup pairing on $H^1(\OO_T)$ is perfect alternating \citep[Proposition~7.27 and proof, p.~108; Corollary~6.13, pp.~93--94]{EGM}. Tensor it with $J_2$. By Serre duality, its restriction to $S$ is induced by $H^1(aJ_2a^\vee)$, which is zero because it factors through the acyclic $\mathfrak g$.

The classes $\alpha,\zeta$ are independent: $F\alpha=\alpha\ne0$, while $F\zeta=0$ and $\zeta\ne0$. Connectedness proves the former nontriviality; a trivial $\alpha _2$-torsor would have nonsmooth total space, proving the latter. They form a basis of $H^1(\OO_T)$. Since $\zeta\cup-$ is invertible on the pencil, $S$ is the graph
$\{\alpha\otimes v+\zeta\otimes T_sv\}$. Isotropy gives
$J_2T_s+T_s^tJ_2=0$. In dimension two this is
$(\operatorname{tr}T_s)J_2=0$.

We verify that this is the actual trace-lift matrix used for Cartier, rather than an unrelated cup matrix. A trace lift $q_i$ can locally be written $z_js_i+r_{i,j}$, with trace-one Artin--Schreier generators $z_j$. Its contraction descends, and the local section $r_{i,j}s_i$ lifts that contraction in (C6). The differences are the torsor cocycle times $a_i$, so its boundary is $\alpha\cup a_i$. The norm sections have boundaries $\zeta\cup a_i$. Projection along the Cartier image into the norm-section basis therefore gives precisely $T_s$, with the relative Frobenius identifications. The derivative comparison proved below transfers trace zero to the actual variations.

Finally, if a nonzero trace-zero singular matrix occurs, conjugate it to $N=\left(\begin{smallmatrix}0&1\\0&0\end{smallmatrix}\right)$. Another trace-zero matrix is $B=\left(\begin{smallmatrix}a&b\\c&a\end{smallmatrix}\right)$. If the entire linear space is singular, then
$\det(N+tB)=tc+t^2(a^2+bc)$ vanishes identically. Thus $c=a=0$ and $B$ is a multiple of $N$. This proves the alternative.
\end{proof}

\begin{lemma}[Rank-one matrix pencils]\label{lem:matrix-pencil}
Over an infinite field, a linear subspace of two-by-two matrices consists entirely of singular matrices if and only if it has a common nonzero kernel vector or a common image line.
\end{lemma}
\begin{proof}
Either common subspace forces singularity. Conversely take a nonzero member and change its source and target bases to make it $\operatorname{diag}(1,0)$. For another member $(b_{ij})$, the polynomial determinant along the line joining them gives $b_{22}=0$, and its own determinant gives $b_{12}b_{21}=0$. Thus the matrix space is the union of the two subspaces defined by $b_{12}=0$ and $b_{21}=0$. A vector space cannot be the union of two proper linear subspaces. In the first case all matrices kill the second basis vector; in the second all images lie in the first coordinate line. The zero space satisfies both alternatives.
\end{proof}

\begin{lemma}[The Frobenius-typed derivative]\label{lem:derivative}
Let $K$ and $Q$ be the kernel and cokernel of the restricted coherent Frobenius at the pullback section of Lemma~\ref{lem:double}. Use square-root coordinates $u$ on the source section space, with ordinary coordinates $u^{[2]}$. In compatible relative Frobenius bases the trace-lift matrix is a rational matrix $A(u)$ of trace zero. Its root-removed Cartier-dual variation is $dA_u(v):Q^*\to K^*$. The coherent variation $K\to Q$ in the actual section direction $v^{[2]}$ has matrix
\[
                            (dA_u(v)^{[2]})^t.            \tag{F3}
\]
\end{lemma}
\begin{proof}
Let $U$ be the root-coordinate space and $W=U^{(2)}$. Root-removed Cartier multiplication is an alternating symmetric bilinear map $B:U\otimes U\to W$, zero between different weights. Set $M(u)v=B(u,v)$,
$r(u)=(u_0,u_1)$ and $R(u)=(u_0^{[2]},u_1^{[2]})$.
Sharpness gives $W=\operatorname{im}M(u)\oplus\operatorname{im}R(u)$.
Let $P(u)$ be projection onto the second summand in its displayed basis:
\[
                         PM=0,\qquad PR=1.               \tag{F4}
\]
Choose a linear weight-preserving trace splitting $q:U\to\widetilde U$, and its twist $q^{(2)}$ on $W$. There are bilinear maps $H,J,K$ defined by
\[
 \widetilde B(\pi^*u,qv)=\pi^*H(u,v)+q^{(2)}B(u,v),\quad
 \widetilde B(qu,qv)=\pi^*J(u,v)+q^{(2)}K(u,v).
\]
The deck involution sends $qv$ to $qv+\pi^*v$. Expanding equivariance gives
$K(u,v)=H(u,v)+H(v,u)+B(u,v)$.
The trace-lift matrix is $A=PH(u)r(u)$. Replacing $q$ by $q+\pi^*\Lambda$ adds
$B(u,\Lambda v)+\Lambda^{[2]}B(u,v)$ to $H(u,v)$; on $v=u_i$ the first term is killed by $P$ and the second is zero. Thus $A$ is independent of this splitting.

Choose correction columns $w_i$ with
$H(u,u_i)+M(u)w_i=R(u)A_i$. The corrected lift is $qu_i+\pi^*w_i$. Its Cartier-dual trace-direction variation has column
$P(K(v,u_i)+B(v,w_i))$, by the trace-pairing formula in Lemma~\ref{lem:double}. Differentiate (F4). Since $dR=0$, one obtains
$(dP)R=0$ and $(dP)M=P\,dM$. Differentiating $A$ and substituting the correction equation gives
\[
 dA(v)_i=P B(v,w_i)+P H(v,u_i)+P H(u,v_i).
\]
Weight orthogonality identifies the last term with $PH(u_i,v)$; the remaining term $PB(v,u_i)=PB(u,v_i)$ vanishes. The deck identity therefore gives exactly the preceding variation column. Coefficient Frobenius and dualization give (F3). The cup comparison in Theorem~\ref{thm:tracezero} gives $\operatorname{tr}A=0$, so its differential and both transported variations have trace zero.

On a constant-rank chart of the sharp section open, the projection and corrected trace lifts can be chosen regularly by inverting a nonzero minor. The trace identities just proved on the dense open of disjoint zero schemes are therefore identities of regular matrices on that chart. They extend to every sharp section on it. This explains the scope of the first-variation assertion even where the determinant-curve construction itself is not a bundle construction.
\end{proof}

\section{The actual determinant directions and the residual boundary}
\label{sec:boundary}

This section identifies the actual determinant directions with a restricted cup product and computes its full residual boundary kernel.

All the bundles below live on the smooth surface $T$. No smoothness assumption on $D_s$ is imposed unless stated. Put $U_s=H^1(Q_s^\vee)$, a two-dimensional space at a sharp section, and define the actual determinant directions by
\[
 \Delta_s=\operatorname{im}\bigl[H^0(E_0)\oplus H^0(E_1)
 \longrightarrow H^0(\mathcal L)\bigr],\qquad
 (t_0,t_1)\longmapsto s_0\wedge t_1+t_0\wedge s_1.
\]

\begin{proposition}\label{prop:restricted}
The derivative of the cohomology subspace $S$ in (F2) is
\[
 \Psi_s(v)=H^1(vJ_{Q_s}):U_s\longrightarrow H^1(Q_s)
       \quad(v\in\Delta_s),                               \tag{P1}
\]
where $J_{Q_s}:Q_s^\vee\to Q_s\mathcal L^{-1}$ is the inverse alternating identification. The determinant-direction sequence is
\[
 0\to H^0(K_sP)\to H^0(E_0)\oplus H^0(E_1)\to
 H^0(\mathcal L)\to H^1(K_sP)\to0,                          \tag{P2}
\]
so $\operatorname{codim}\Delta_s=h^1(K_sP)=h^0(K_sP)$.
Define the symmetric cup-wedge map
\[
 \mu_s:\operatorname{Sym}^2U_s\longrightarrow H^2(\mathcal L^{-1}),
 \qquad xy\longmapsto\wedge_*(x\cup y).
\]
Under Serre duality and restriction to $\Delta_s^*$, this is the dual of $\Psi_s$. In particular an invertible variation exists at this sharp source if and only if
\[
 \{x\in U_s:\mu_s(xy)|_{\Delta_s}=0\text{ for every }y\in U_s\}=0.
                                                                    \tag{P3}
\]
\end{proposition}
\begin{proof}
The middle bundle in (F1) is fixed as the pair varies. For a variation of its Lagrangian inclusion $i$, the derivative of the subspace is represented by $p\dot i$, where $p$ is the quotient map. A frame change adds $iX$, killed by $p$. Locally normalize $s_0=(1,0)$ using a determinant-one frame even over dual numbers. The infinitesimal action of this frame change on $\operatorname{End}^0(V)/\OO$ is zero: the commutator of two trace-zero rank-two matrices is scalar in characteristic two. Thus the frame of $\mathfrak g$ is stationary to first order. In (C4), $d(h^2)=0$ and $u$ is parameter-independent, so
\[
 \dot a=\begin{pmatrix}0&0\\0&\dot w\end{pmatrix},\qquad
 \dot f=\begin{pmatrix}u\dot w&0\\0&0\end{pmatrix}.
\]
It follows that $p\dot i=aJ_2\dot a^t+f(u^{-1}J_2)\dot f^t=\dot wJ_2$.
This is the intrinsic map $\dot wJ_{Q_s}$. Differentiating a \v Cech representative of $S$ adds terms inside $S$ and the cohomology class of this map; modulo $S$ only the latter survives. Cohomology dimensions are constant on the sharp open, so this computation proves (P1), equivalently the derivative of the actual graph matrix $T_s$.

Twist the contraction sequence defining $K_s$ by $P$, then exchange the two summands. Its map is now the determinant-direction map. It is surjective as a sheaf because the sections have no common zero. Higher cohomology of the $E_i$ vanishes, giving (P2). Both middle section dimensions equal $m$, so its Euler characteristic gives the codimension identity. The twist $P$ is part of this exact sequence and cannot be omitted.

Interchanging two degree-one cohomology factors and then the two arguments of the alternating bundle pairing contributes two signs. Hence $\mu_s$ is symmetric, also in characteristic two; no identification with divided powers is used. For $x,y\in U_s$ and $v\in\Delta_s$, Serre duality gives
\[
 \langle H^1(vJ_{Q_s})x,y\rangle
       =\operatorname{Tr}_T\bigl(v\,\wedge_*(x\cup y)\bigr).
\]
This proves the asserted duality. The radical in (P3) is the common kernel of the variation family. An invertible member has zero common kernel. Conversely Theorem~\ref{thm:tracezero} says that an entirely singular family is zero or a single nilpotent line, which has a nonzero common kernel. This proves the equivalence.
\end{proof}

\begin{theorem}[Residual-pencil boundary and its full kernel]\label{thm:boundary}
Put $E=Q_s^\vee$, $G=\mathfrak g^\vee$, and $R=R_s$. There are exact sequences
\[
 0\to E\xrightarrow{\varphi}G\to i_*R\to0,                  \tag{P4}
\]
\[
 0\to\mathcal L^{-1}\xrightarrow{d_2}E\otimes G
 \xrightarrow{d_1}\operatorname{Sym}^2G\to i_*R^2\to0,       \tag{P5}
\]
with $d_1(e\otimes g)=\varphi(e)g$ and
$d_2(e\wedge e')=e\otimes\varphi(e')-e'\otimes\varphi(e)$.
The boundary in (P4) identifies $H^0(R)=U_s$. Under this identification,
\[
 \mu_s=\partial_s\circ\operatorname{mult}:
 \operatorname{Sym}^2H^0(R)\to H^0(R^2)\to H^2(\mathcal L^{-1}), \tag{P6}
\]
where $\partial_s$ is the two-step boundary of (P5). Its full kernel fits into
\[
\begin{split}
 0\to H^0(\operatorname{Sym}^2G)\to\ker\partial_s
 \to\ker[H^1(E\otimes G)\to H^1(\operatorname{Sym}^2G)]\to0,\\
 H^1(E\otimes G)=H^0(D_s,A_s\otimes G|_{D_s}).
\end{split}\tag{P7}
\]
The last domain is nonzero. If $D_s$ is integral, the residual products form a three-dimensional subspace $W_s\subset H^0(R^2)$, and
\[
 \dim\operatorname{im}\Psi_s
 =\dim\operatorname{im}\bigl[W_s\xrightarrow{\partial_s}
       H^0(\mathcal L)^*\longrightarrow\Delta_s^*\bigr].     \tag{P8}
\]
\end{theorem}
\begin{proof}
Dualize $0\to\mathfrak g\to Q_s\to i_*A_s\to0$ from Lemma~\ref{lem:factor}. Cartier-divisor duality gives (P4), since
$\mathcal E xt^1(i_*A_s,\OO_T)=i_*(A_s^{-1}\mathcal L|_{D_s})=i_*R$.
Acyclicity of $G$ gives its boundary isomorphism. Locally the map $\varphi$ has diagonal form $(w,1)$ after absorbing a unit. On bases $e_i,f_i$ its four tensor images are
$wf_1^2,wf_1f_2,f_1f_2,f_2^2$. Their only relation is
$e_1\otimes f_2-w e_2\otimes f_1$, because $w$ is a nonzerodivisor. The quotient is $(\OO/(w))f_1^2$, proving (P5) even for nonreduced $D_s$.

To verify (P6) on both off-diagonal and diagonal products, lift $r,s\in H^0(R)$ locally to $u_i,v_i\in G$. Write
$u_j-u_i=\varphi(x_{ij})$ and $v_j-v_i=\varphi(y_{ij})$.
A lift under $d_1$ of the difference of $u_iv_i$ is
$z_{ij}=x_{ij}\otimes v_j+y_{ij}\otimes u_i$.
Its next difference is
$-d_2(x_{ij}\wedge y_{jk})$.
In characteristic two this is the cup-wedge cocycle with positive sign. The same calculation works for $r=s$, proving (P6) without division by two.

The dual of (C1), tensored by $G$, is
$0\to E\otimes G\to G^2\to i_*(A_sG|_{D_s})\to0$.
Acyclicity gives the second line of (P7) and $H^0(E\otimes G)=0$.
Let $\mathcal J$ be the image of $d_1$. Since
$H^1(\mathcal L^{-1})=0$, the first short sequence in (P5) gives
$H^0(\mathcal J)=0$ and an injection
$H^1(E\otimes G)\to H^1(\mathcal J)$.
The second short sequence gives the boundary
$H^0(R^2)\to H^1(\mathcal J)$ with kernel
$H^0(\operatorname{Sym}^2G)$ and image the kernel of
$H^1(\mathcal J)\to H^1(\operatorname{Sym}^2G)$.
The inverse image of $H^1(E\otimes G)$ is precisely $\ker\partial_s$, proving the first line of (P7).

The bundle $G$ has a subline $P$, as proved explicitly in the next lemma. Curve duality and $\chi(A_sP)=0$ give
$h^0(A_sP)=h^0(RP)\ge1$, because the norm section $\sigma$ is nonzero. Thus $H^0(A_sG|_{D_s})\ne0$. Finally, on an integral curve the ratio of two independent sections of $R$ is nonconstant. A relation between their three quadratic monomials would make that ratio algebraic over the algebraically closed field, hence constant. Multiplication is therefore injective on its three-dimensional symmetric square. Combining (P6) with Proposition~\ref{prop:restricted} proves (P8).
\end{proof}

\begin{lemma}[The unipotent bundle and its Frobenius]\label{lem:unipotent}
Let $b(H)$ denote the local index of the actual section bundle $H$. The bundle
$U=\operatorname{End}^0(V)/\OO_T$ is an extension
$0\to\OO_T\to U\to\OO_T\to0$ with class $\gamma$ satisfying $F\gamma=0$.
It is nonsplit if $b(H)=0$, and split if $b(H)>0$.
Moreover $G\simeq UP$ and
\[
 \begin{array}{c|c|c}
 &\operatorname{Sym}^2G&(h^0,h^1,h^2)\\ \hline
 b(H)=0&U\oplus\OO_T&(2,4,2)\\
 b(H)>0&\OO_T^3&(3,6,3).
 \end{array}                                               \tag{P9}
\]
\end{lemma}
\begin{proof}
Conjugation on trace-zero rank-two matrices modulo scalars is the Frobenius-twisted standard representation divided by the determinant. Thus $U=F_{\mathrm{abs},T}^*V\otimes L^{-1}$, interpreted semilinearly, or in relative notation $F_T^*(V^{(2)})L^{-1}$.
Choose $\alpha _2$-torsor charts $z_i^2=h_i$, $z_j=z_i+c_{ij}$, and transitions $a+bz_i$ for $J$. The transition matrix of $V$ is
\[
 M_{ij}=\begin{pmatrix}a&ac_{ij}+bh_i\\b&bc_{ij}+a\end{pmatrix},
 \qquad N_{ij}=a^2+b^2h_i.
\]
The normalized square $M_{ij}^{[2]}/N_{ij}$ sends $(h_j,1)^t$ to $(h_i,1)^t$. In frames with second column $(1,0)^t$ it becomes
$\left(\begin{smallmatrix}1&\gamma_{ij}\\0&1\end{smallmatrix}\right)$, where
\[
                         \gamma_{ij}=b^2/N_{ij}.           \tag{P10}
\]
This constructs the actual subline and quotient, not just their pullbacks.

If $\delta$ is the invariant vector field with $\delta(z_i)=1$, the cocycle
$\kappa_{ij}=b/(a+bz_i)$ represents $d\lambda_J(\delta)$: infinitesimal translation changes the transition by $1+\epsilon\kappa_{ij}$. Let $\nu^c:T\to C^{(2)}$ be the complementary isogeny, so $\nu^c\nu=F_C$. Its kernel is $\mu _2$, and its dual kernel is \'etale. Hence pullback by $\nu^c$ is an isomorphism on $H^1(\OO)$, by the differential of the dual isogeny \citep[Theorem~7.5 and proof, p.~100]{EGM}. Direct substitution gives
$(\nu^c)^*\kappa^{(2)}=\gamma$.
The local-local block of $\lambda_J$ is an isomorphism when $b(H)=0$, so its differential does not kill the $\alpha _2$ tangent. If $b(H)>0$, that block is divisible by $[2]$ and kills it. The local index of $J$ equals that of $H$, since $\rho$ is an isomorphism on the local-local factor. This proves exactly the split and nonsplit assertions.

The forms $dh_i$ glue to a nonzero invariant form $\theta$: smoothness of the torsor's total space forces $dh_i\ne0$. Differentiating (P10) gives $d\gamma_{ij}=\gamma_{ij}^2\theta$. Hodge--de Rham degeneration for abelian varieties makes $d:H^1(\OO_T)\to H^1(\Omega_T^1)$ zero \citep[Corollary~7.29 and proof, pp.~108--109]{EGM}. Multiplication by the nonzero invariant $\theta$ injects $H^1(\OO_T)$ into that group, so $F\gamma=0$.

Write $\gamma_{ij}^2=f_j-f_i$. On the symmetric-square frame the transition of $U$ is
$\left(\begin{smallmatrix}1&\gamma_{ij}&\gamma_{ij}^2\\0&1&0\\0&0&1\end{smallmatrix}\right)$.
Replacing its third basis vector locally by itself plus $f_i$ times the first removes the last off-diagonal term. Thus $\operatorname{Sym}^2U=U\oplus\OO$; $U$ is self-dual and $P^2=\OO$, giving the same assertion for $G=UP$. For nonzero $\gamma$, cup product with it has rank one in degrees zero and one, because the $H^1(\OO)$ pairing is perfect alternating. The extension sequence gives $h^\bullet(U)=(1,2,1)$, hence (P9). The split case is immediate. In particular an even exponent can change (P9), even when the root polarization has local index zero.
\end{proof}


\begin{proposition}[Automorphisms and the direction codimension]\label{prop:auto-directions}
For the general sharp pair above there is a natural exact sequence
\[
 0\longrightarrow\End^0(V)\xrightarrow{\Theta_s}K_sP
   \longrightarrow i_*(A_sP|_{D_s})\longrightarrow0,
 \qquad \Theta_s(u)=(us_1,us_0).                         \tag{P11}
\]
Let $\delta_{\mathrm{auto}}:H^0(D_s,A_sP|_{D_s})
\to H^1(T,\End^0(V))$ be its boundary. Then
\[
 \operatorname{codim}_{H^0(\mathcal L)}\Delta_s
 =\epsilon_b+\dim\ker\delta_{\mathrm{auto}},\qquad
 \epsilon_b=h^0(\End^0(V))=
 \begin{cases}1,&b(H)=0,\\2,&b(H)>0.\end{cases}            \tag{P12}
\]
In particular positive local index forces direction codimension at least two.
No nonzero kernel of $\delta_{\mathrm{auto}}$ is asserted.
\end{proposition}
\begin{proof}
For a rank-two alternating pairing, direct multiplication gives
\[
 \beta(s_0,us_1)+\beta(s_1,us_0)
       =\operatorname{tr}(u)\beta(s_0,s_1).
\]
Using the fixed $P^2\simeq\OO_T$, this shows that $\Theta_s$ takes values in
$K_sP$. The sections $s_0,s_1$ form a basis at the generic point after a
rational frame of $P$, so $\Theta_s$ is injective. Scalars map to the
crossing section $(s_1,s_0)$. Consequently it defines a commutative diagram
of the two extensions
\[
\begin{array}{ccccccccc}
0&\to&\OO_T&\to&\End^0(V)&\to&\End^0(V)/\OO_T&\to&0\\
 &&\Vert&&\downarrow\Theta_s&&\downarrow f\otimes P\\
0&\to&\OO_T&\to&K_sP&\to&Q_sP&\to&0.
\end{array}
\]
The right-hand map is the map of Lemma~\ref{lem:factor} tensored by $P$;
its kernel is zero and its cokernel is $i_*(A_sP|_{D_s})$. The snake lemma
proves (P11), including its global cokernel line.

Here is a local verification that retains the frame issue in characteristic
two. If $p$ frames $P$ with $p^2=u\in\OO_T^*$, first pass to the finite
faithfully flat cover defined by $t^2=u$. Its invertible $t$ makes $p/t$
a square-normalized frame; no Zariski square root is presumed. On a chart
where $s_0$ is nonzero, use a determinant-one frame with $s_0=(1,0)$ and
write $s_1=(a,b)$. A trace-zero endomorphism has matrix
$\left(\begin{smallmatrix}r&x\\y&r\end{smallmatrix}\right)$.
Write a pair in the contraction kernel as
$((u_0,aw+bv),(v,w))$. Then
\[
 (r,x,y)\longmapsto(ar+bx,r,y),\qquad
 \begin{pmatrix}a&b&0\\1&0&0\\0&0&1\end{pmatrix}
\]
has determinant $b$. Its scalar direction is $(a,1,0)$; quotienting by
that direction gives $(x,y)\mapsto(bx,y)$, exactly $f\otimes P$.
The chart where $s_1$ is nonzero is symmetric. These intrinsic maps and
their exactness descend by faithful flatness.

We compute the global trace-zero endomorphisms, keeping the local index of
$H$. In the torsor coordinates of Lemma~\ref{lem:unipotent}, pullback to
$C$ gives
\[
 \nu^*V=J\otimes U_\kappa,
 \qquad 0\to\OO_C\to U_\kappa\to\OO_C\to0.
\]
Indeed, on $C\times_T C$ write the second torsor coordinate as $z+e$, with
$e^2=0$. A transition $g=a+bz$ for $J$ becomes
$g(z+e)/g(z)=1+e\,b/(a+bz)$. Thus its extension cocycle is precisely
$\kappa_{ij}=b/(a+bz_i)=d\lambda_J(\delta)$, where
$\delta(z_i)=1$. Lemma~\ref{lem:unipotent} proves that this class is nonzero
exactly when $b(H)=0$.
Adjunction now gives
\[
 \operatorname{Hom}_T(V,V)
 =\operatorname{Hom}_C(\nu^*V,J)=H^0(C,U_\kappa^\vee).
\]
When $\kappa\ne0$, the extension boundary from $H^0(\OO_C)$ is nonzero,
so this space is one dimensional. All global endomorphisms are scalar,
and their trace is zero in rank two and characteristic two.

When $\kappa=0$, the last displayed space has dimension two. The vanishing
of the cocycle also glues local differentiation along $\delta$ to a global
operator $D:J\to J$ satisfying
$D(fv)=\delta(f)v+fD(v)$. This follows directly by correcting the local
operators by functions whose differences are $\delta(g)/g=\kappa_{ij}$.
The operator is $\OO_T$-linear because $\delta$ kills functions from $T$.
Locally, in a frame $v$ for $J$, write $D(v)=(a+bz)v$ and $z^2=h$.
On the $\OO_T$-basis $v,zv$ its matrix is
\[
                 \begin{pmatrix}a&1+bh\\b&a\end{pmatrix}.
\]
Its trace is zero. It is independent of the identity because its Leibniz
rule has the nonzero derivative $\delta(z)=1$. These two operators span
all global endomorphisms, so all of them are traceless. This proves the
formula for $\epsilon_b$.

Taking global sections in (P11) yields
\[
 h^0(K_sP)=h^0(\End^0(V))+\dim\ker\delta_{\mathrm{auto}}.
\]
The determinant-direction sequence (P2) has source and target of the same
dimension $m$, so $\operatorname{codim}\Delta_s=h^0(K_sP)$.
Together these identities prove (P12). Although the norm argument gives
$H^0(A_sP|_{D_s})\ne0$, its boundary may be injective. It supplies no
additional positive term in (P12) without control of that boundary.
\end{proof}

\section{Smooth determinant curves and the exact open question}
\label{sec:open}

This section establishes smoothness under explicit jet hypotheses and states the remaining smaller-degree and determinant questions.

\begin{proposition}\label{prop:jets}
Suppose $J$ and $JQ_0$ are cubes of very ample bundles. Then $E_0,E_1$ separate first jets on $T$. A general pair has disjoint reduced zero schemes and smooth integral determinant curve $D_s$. Its squared pencil is separable. If a sharp open is nonempty, these conditions can hold there, and the ordinary curve multiplication
\[
 H^0(A_s)\otimes H^0(\omega_{D_s}A_s^{-1})
          \longrightarrow H^0(\omega_{D_s})
\]
has rank three or four, with rank three exactly when
$\omega_{D_s}A_s^{-2}$ is trivial. The power hypothesis holds, for example, when $M$ itself descends through $\rho$ and $H=M^{3e}$ with $e\ge3$.
\end{proposition}
\begin{proof}
Completed regular parameters put the $\alpha _2$ quotient locally in the form
$k[[x,z^2]]\to k[[x,z]]$. The inverse image of a first infinitesimal neighbourhood has ideal $(x^2,xz^2,z^4)$ and basis
$1,z,z^2,z^3,x,xz$. This ideal contains $(x,z)^4$.
The cube of a very ample bundle separates third jets: take a nonvanishing section and sections whose ratios generate the cotangent space, then take products of three to span all monomials of degree at most three. Thus both finite pushforwards separate first jets.

For a rank-two section, vanishing imposes two conditions, and failure of its derivative to have rank two adds one. The resulting incidence has codimension three over a surface, so a general zero scheme is reduced and zero-dimensional. Simultaneous vanishing of the pair imposes four conditions. On its complement the determinant matrix has rank one. Normalize a nonzero value to the first basis vector; determinant zero and its two derivative equations impose three independent jet conditions on the other section. Thus its singular-point incidence also has codimension three, and a general determinant curve is smooth. It is connected by the ample divisor sequence, hence integral.

If the image of $\eta_P$ in $\Omega_{D_s}^1$ were zero, the conormal sequence would put its restriction in $H^0(\OO_{D_s}(-D_s))=0$. But a nonzero invariant form on $T$ is nowhere vanishing. Therefore (C8) proves separability. These conditions are nonempty opens and meet any nonempty sharp open.

At sharpness both pencil section spaces have dimension two. The base-point-free pencil sequence identifies the kernel of the displayed multiplication with $H^0(\omega_{D_s}A_s^{-2})$. Its degree is zero on the smooth integral curve, so its dimension is one exactly in the trivial case, and zero otherwise. Finally, if $J=J_0^{3e}$ with $e\ge3$, then $J_0^e$ is very ample. A cube root of the degree-zero twist $Q_0$ is a translate after tensoring with $J_0^e$, since the polarization map is surjective on geometric points. This verifies the stated sufficient power condition.
\end{proof}

The primitive rank question is
\[
       r_e=9\chi(M)e^2-2\quad\text{for }a(M)>0,\ b(M)=0.
\]
Theorem~\ref{thm:eventual-all-mixed} proves this equality for every sufficiently large $e$ on every mixed pair, including surfaces with no finite-field model. Theorem~\ref{thm:ample-difference} gives an explicit range when an ordinary-index-zero ample bundle exists, and Theorem~\ref{thm:finite-field} gives the sufficient bound~\eqref{eq:finite-threshold} for a descended pair. Theorem~\ref{thm:no-ordinary-index-zero-ample-class} shows why the auxiliary class cannot be assumed on every surface: its example has ordinary index $1+v_2(t)$ on every ample class $t[L]$. Eventual rank nevertheless holds there. The remaining rank questions concern the smaller degrees outside the exact ranges of Theorems~\ref{thm:p10exact}--\ref{thm:p10bounds} and quantitative bounds for arbitrary pairs. No uniform threshold across all pairs is established.

The sharp first-variation mechanism would establish successive positive indices if the common radical in (P3) vanished at a suitable sharp source. The equivalent residual formulation is (P8), with the actual direction space $\Delta_s$. Neither the existence of a smooth separable pencil nor its ordinary rank-three multiplication proves that condition. The relevant products are $\operatorname{Sym}^2H^0(R_s)$, followed by the two-step boundary and then restriction to $\Delta_s$. Formula (P7) shows why the acyclicity of $G$ does not remove the boundary kernel: a further twisted-pencil domain is already nonzero. Likewise the dimensions in (P9) determine only one part of that kernel. Formula (P12) identifies the unavoidable automorphism contribution to the direction codimension, but neither proves that $\ker\delta_{\mathrm{auto}}$ is nonzero nor removes the restriction to $\Delta_s$. Uniform nonvanishing of this actual first variation remains open. The ample-difference construction and the spreading theorem establish their rank conclusions without settling this particular determinant test.

\appendix
\section{Cohomology, stable rank, and nodal curves}\label{app:cohomology}

This appendix relates stable coherent rank to curve defects and explains its behavior under nodal degeneration. We use the usual polarization homomorphism
\[
 \lambda_N:A\longrightarrow A^\vee=\Pic^0(A),\qquad
 a\longmapsto t_a^*N\otimes N^{-1}.
\]
For an ample $N$, $\lambda_N$ is an isogeny, $\deg\lambda_N=\chi(N)^2$, and $\chi(N)=N^2/2$ on a surface. A numerical equivalence of line bundles is denoted by $\equiv$. On an abelian variety it is equivalence modulo $\Pic^0$, and $\NS(A)$ is torsion-free. Indeed $[N]\mapsto\lambda_N$ embeds it into the torsion-free group $\Hom(A,A^\vee)$; see \citet[\S7 and Proposition~11.2--Corollary~11.3, pp.~159--160]{EGM}.

\begin{lemma}\label{lem:cohomology}
If $N$ is ample on an abelian surface, then $H^i(A,N)=0$ for $i>0$ and $H^1(A,N^{-1})=0$. For every effective divisor $C\in|N|$, the curve is connected, its arithmetic genus is $\chi(N)+1$, and there is a Frobenius-compatible exact sequence
\begin{equation}\label{eq:relative}
 0\longrightarrow H^1(A,\OO_A)\longrightarrow H^1(C,\OO_C)
 \longrightarrow H^2(A,N^{-1})\longrightarrow H^2(A,\OO_A)
 \longrightarrow0.
\end{equation}
On $H^2(A,N^{-1})$, the semilinear operator in this sequence is the composite
\[
 H^2(A,N^{-1})\xrightarrow{F}H^2(A,N^{-2})\xrightarrow{\,s\,}H^2(A,N^{-1}).
\]
The last map of the sequence is dual to the inclusion $k s\subset H^0(A,N)$, where $s$ defines $C$. In particular its kernel has dimension $\chi(N)-1$.
\end{lemma}
\begin{proof}
An ample bundle has a symmetric translate. To see this, write $[-1]^*N\otimes N^{-1}=P\in\Pic^0(A)$, choose $Q$ with $Q^2=P$, and observe that $N\otimes Q$ is symmetric. Surjectivity of $\lambda_N$ realizes this twist as a translate. For a symmetric representative and an arbitrarily large odd integer $u$, one has $[u]^*N=N^{u^2}$. The trace of the finite \'etale map $[u]$, whose degree is the invertible scalar $u^4$, splits the map on structure sheaves. Projection formula therefore injects $H^i(A,N)$ into $H^i(A,N^{u^2})$. Serre vanishing proves the first assertion; duality and $\omega_A\simeq\OO_A$ prove the second. Translation gives the assertion for the original bundle and all its algebraically trivial twists.

Apply cohomology to $0\to N^{-1}\xrightarrow{s}\OO_A\to\OO_C\to0$. Since $H^0(N^{-1})=H^1(N^{-1})=0$, one obtains $H^0(\OO_C)=k$ and the displayed sequence. Its terminal surjectivity follows either from $H^2(C,\OO_C)=0$ or from duality. The dual map sends $1$ to $s$. Riemann--Roch gives $h^0(N)=\chi(N)$ and adjunction gives $p_a(C)=1+N^2/2$. Frobenius compatibility follows from the restriction-sequence map whose left-hand operation sends a local section $u$ of $N^{-1}$ to $s u^2$: multiplication by the defining section then gives $s^2u^2=(su)^2$ in $\OO_A$.
\end{proof}

\begin{lemma}\label{lem:stable}
Stable parts are exact on finite-dimensional semilinear Frobenius spaces over $k$. In a proper flat family of connected divisors from a fixed complete linear system on $A$, the locus where the stable rank of $H^1(\OO_C)$ is at least a specified integer is open. The same assertion holds for the relative quotient on smooth fibres.
\end{lemma}
\begin{proof}
For a space of dimension $b$, the images of iterated Frobenius stabilize after at most $b$ steps. Since $k$ is perfect, their dimensions behave exactly as for linear maps. The stabilized image is bijective under Frobenius and its complementary generalized kernel is nilpotent; the direct-sum decomposition follows by the usual Fitting argument applied to a sufficiently high iterate. A map respects both parts. Given a surjection, lift a vector in the stable part and repeatedly apply Frobenius after inverting it on the stable target; the nilpotent components disappear, proving surjectivity on stable parts. Injectivity and exactness in the middle follow from the same decomposition.

The preceding lemma makes $h^0(\OO_C)=1$ and $h^1(\OO_C)=\chi(N)+1$ constant throughout the divisor family. Cohomology and base change give a vector bundle carrying a linearized Frobenius map. In a local basis with matrix $B$, the $b$th iterate has matrix $B B^{[2]}\cdots B^{[2^{b-1}]}$, with the order depending only on the chosen linearization convention. Its rank is the stable rank on a geometric fibre. Nonvanishing of a minor is open. The relative spaces are the constant-rank kernels in \eqref{eq:relative}, so the same argument applies there. Exactness of stable parts gives \eqref{eq:dimensions}, using the standard cohomological description of the $2$-rank.
\end{proof}

\begin{lemma}[Defect of a nodal curve]\label{lem:nodal}
Let $C$ be a connected reduced projective nodal curve with smooth normalization components $C_1,\ldots,C_v$. If its dual graph has first Betti number $b_1$, then
\[
 h^1(C,\OO_C)=\sum_i g(C_i)+b_1,\qquad
 \dim H^1(C,\OO_C)_{\mathrm{st}}=\sum_i f_{C_i}+b_1.
\]
Thus its defect is the sum $\sum_i(g(C_i)-f_{C_i})$.
\end{lemma}
\begin{proof}
The normalization sequence is
\[
 0\to\OO_C\to\bigoplus_i\nu_*\OO_{C_i}\to\bigoplus_{x\text{ node}}k_x\to0,
\]
where the last map is the difference of the two branch values. The induced map on constants is the graph coboundary. Its cokernel is $k^{b_1}$ and gives an exact sequence
$0\to k^{b_1}\to H^1(C,\OO_C)\to\bigoplus_iH^1(C_i,\OO_{C_i})\to0$.
Frobenius is coordinatewise squaring on the graph term and hence is bijective. Apply Lemma~\ref{lem:stable}.
\end{proof}

We will construct nodal members with a given defect and use openness to produce smooth members of no larger defect. Very ampleness supplies a nonempty smooth-member open; irreducibility of the complete linear system makes these two opens meet.

\section{Theta divisors and jets in characteristic two}
This appendix supplies the point and jet separation needed to control the singularities of finite norm images.


\begin{lemma}\label{lem:dichotomy}
A principally polarized abelian surface has either a product theta divisor on a product of elliptic curves or a smooth genus-two theta divisor. In the second case the polarized surface is the Jacobian of that curve.
\end{lemma}
\begin{proof}
Write an effective theta divisor as $\Theta=\sum m_iB_i$. Abelian varieties contain no rational curves: after removing any inseparable part from a putative map from $\PP^1$, a nonzero differential would pull invariant differentials back to nonzero regular differentials on $\PP^1$. For every integral $B_i$, adjunction therefore gives $B_i^2\ge0$, with this integer even. Since $\Theta^2=2$ and every $\Theta B_i$ is positive, reducibility forces precisely two reduced components. A divisor $mB$ with $m>1$ is impossible because $m^2B^2=2$. Their equations $B_i^2+B_1B_2=1$ imply $B_i^2=0$ and $B_1B_2=1$. They are smooth elliptic translates, and addition of the corresponding elliptic subgroups is an isogeny of degree one. This identifies the polarization with the product polarization.

If $\Theta$ is integral, its arithmetic genus is two. Its normalization cannot have genus one, since an elliptic curve maps to an elliptic translate, of self-intersection zero, on an abelian variety. Thus $\Theta$ is smooth of genus two. Its inclusion induces a surjection $u:J(\Theta)\to A$ and hence an isogeny. If $\Theta_J$ is the Abel--Jacobi divisor, then $(u^*\Theta)\Theta_J=2$ and $\Theta_J^2=2$. Hodge index gives $2\deg u=(u^*\Theta)^2\le2$, so $u$ is an isomorphism. The standard Jacobian identities used here hold in characteristic two; see \citet[Theorem~14.23, Corollary~14.24, pp.~230--231; Theorem~14.29, p.~233; Example~14.34, pp.~236--237]{EGM}.
\end{proof}

In the smooth-theta case choose a theta characteristic $\kappa$ on the genus-two curve $G$ and embed it as $x\mapsto\OO_G(x)\otimes\kappa^{-1}$. Such a $\kappa$ exists because multiplication by two on $\Pic^0(G)(k)$ is onto. The resulting $\Theta$ is symmetric. Put $P_a=t_a^*\OO(\Theta)\otimes\OO(-\Theta)$ and write $\Theta_a=\Theta+a$.

For a line bundle $L$, at a point $P$, first jets mean evaluation in $L_P\otimes\OO_{A,P}/\mathfrak m_P^2$, and second jets mean evaluation modulo $\mathfrak m_P^3$. Their dimensions on a surface are three and six, respectively.

\begin{lemma}\label{lem:theta-jets}
In the smooth-theta case the following hold.
\begin{enumerate}
\item Every twist of $2\Theta$ is generated by global sections. The system $|2\Theta|$ separates distinct $x,y$ if $y\ne-x$.
\item The system $|3\Theta|$ separates three distinct values and separates second jets at every point. In particular it is very ample.
\item For distinct $P,Q$, the evaluation on their two first jets is onto except possibly when
\begin{equation}\label{eq:badpairs}
 2P+Q\in\Theta\quad\text{and}\quad P+2Q\in\Theta.
\end{equation}
For a fixed nonzero $g$, only finitely many pairs $(P,P+g)$ satisfy \eqref{eq:badpairs}.
\item If $n\ge4$, $|n\Theta|$ separates any two first jets and any three distinct values.
\end{enumerate}
\end{lemma}
\begin{proof}
The divisors $\Theta_c+\Theta_{-c}$ belong to $|2\Theta|$. To be nonzero at a chosen point, avoid the corresponding two translates of $\Theta$. To vanish at $x$ and not at $y\ne\pm x$, take $c$ on $\Theta+x$ outside $(\Theta+y)\cup(\Theta-y)$. The relevant translates are distinct because the stabilizer of a principal theta divisor is trivial. Their pairwise intersections are finite. Every twist of $2\Theta$ is a translate, since $[2]$ is onto on $k$-points. This proves (1); positive-degree cohomology vanishes by Lemma~\ref{lem:cohomology}.

Two elementary facts about $G$ will also be used. Restriction sends $P_a$ to the degree-zero class $-a$ on $G$, and the Gauss map of $\Theta$ is its canonical map $G\to\PP^1$. For the restriction formula, a generic $a=i(r)-i(s)$ gives the intersection $\Theta\cap\Theta_a=\{i(r),i(\iota s)\}$, where $\iota$ is the hyperelliptic involution. Subtracting the canonical divisor $s+\iota s$ proves $i^*P_{-a}=a$. The difference map $G\times G\to A$ is onto because theta translates intersect, and equality extends from the generic case to all $a$. The Gauss assertion follows by restricting invariant one-forms: $H^0(A,\Omega_A^1)\simeq H^0(G,\omega_G)$. The canonical map is finite and onto in characteristic two as well.

For three distinct points $P_1,P_2,P_3$, choose $a\in\Theta+P_1$ and $b\in\Theta+P_2$, with neither in $\Theta+P_3$ and with $P_3+a+b\notin\Theta$. These are simultaneous nonempty open conditions on the product of the two curves: the addition of two theta translates is onto $A$. The product of the sections of $\Theta_a,\Theta_b,\Theta_{-a-b}$ vanishes at $P_1,P_2$ and not at $P_3$. Permuting the points proves three-value separation.

For first jets at $P$, a theta factor through $P$, multiplied by a section of the remaining twist of $2\Theta$ nonzero there, supplies its Gauss linear form. These linear forms span the cotangent space, and a section nonzero at $P$ supplies the value. The same construction separates two distinct points. To obtain second jets, choose $p,q\in\Theta$, put $a=P-p$, $b=P-q$, and multiply the theta sections with parameters $a,b,-a-b$. The first two vanish at $P$, and the third does not vanish when $3P-p-q\notin\Theta$. This is a dense open of $\Theta^2$. Its image under the product of Gauss maps is dense in $(\PP^1)^2$. Products of the resulting pairs of linear forms span all binary quadratics, also in characteristic two. Hence the quadratic terms, as well as the lower jets already obtained, are arbitrary. This proves (2).

For (3), put $s=P+Q$. Theta translates containing $P,Q$ are parametrized by the length-two intersection $Z=(\Theta+P)\cap(\Theta+Q)$. Its support is preserved by $a\mapsto s-a$. A fixed point has $2a=s$ and corresponds to hyperelliptically conjugate points of $G$ with identical Gauss direction; its intersection multiplicity is at least two. Consequently two distinct support points are exchanged, and a unique support point is fixed.

For $a$ in this support, restriction to $\Theta_a$ gives a surjection
\[
 H^0(3\Theta)\longrightarrow H^0\bigl(G,\omega_G^3\otimes(-3a)\bigr).
\]
Write $p=P-a$, $q=Q-a$ on $G$. The degree-six bundle has vanishing $H^1$, and the cokernel of its four tangential jet evaluations is the $H^1$ of the residual degree-two bundle. This has dimension one precisely when that bundle is canonical. Thus the tangential jet conditions have rank four unless
$\omega_G^3\otimes(-3a)(-2p-2q)=\omega_G$, when the rank is three. Since the class of $p+q-K_G$ is $s-2a$, the exception is exactly $a=2s$. The sections containing $\Theta_a$ supply the two normal derivatives: their residual bundle is $2\Theta\otimes P_a$. By (1), these two derivatives are independent unless $a=-s$. Thus the full six jets are obtained whenever a support point avoids both $-s$ and $2s$.

If every support point lies among these two bad parameters, then $-s$ is a support point. Indeed, if the sole support point is $2s$, its fixed-point equation gives $4s=s$ and hence $2s=-s$. Therefore rank failure implies that $\Theta_{-s}$ contains both $P,Q$, which is \eqref{eq:badpairs}. For $Q=P+g$, these conditions are $3P+g\in\Theta$ and $3P+2g\in\Theta$. The two translates are distinct when $g\ne0$, and $[3]$ is finite; their common inverse image is finite.

Finally let $n=3+r$ with $r\ge1$. Choose $(a,b)$ satisfying $3b=-ra$. The parameter space is finite and surjective over either factor, and each component dominates both. The residual factorization is
\[
 \OO(n\Theta)=t_a^*\OO(r\Theta)\otimes t_b^*\OO(3\Theta).
\]
The distinguished section of the first factor is nonzero at prescribed finitely many points on an open in $a$. For a prescribed pair $P,Q$, the bad translated cubic parameters satisfy the two conditions \eqref{eq:badpairs} with $3b$ added; they form a finite set of $b$. Choose a pair $(a,b)$ outside these bad loci. Multiplication by the nonvanishing residual section transfers all six jets. The identical argument transfers the three independent values from (2). This proves (4).
\end{proof}

On a product $E_1\times E_2$, $\OO(3\Theta)=\OO_{E_1}(3O_1)\boxtimes\OO_{E_2}(3O_2)$ is very ample by the plane cubic embeddings. Its Frobenius fibres have length four and are products of length-two elliptic fibres. A degree-three line bundle on an elliptic curve evaluates onto every length-two subscheme, since the kernel has degree one and vanishing $H^1$. We record the consequence also in the smooth-theta case.

\begin{corollary}\label{cor:fgenerated}
For relative Frobenius $F:A\to A^{(2)}$ on a principally polarized surface, $F_*\OO_A(3\Theta)$ is generated by its global sections.
\end{corollary}
\begin{proof}
The local algebra of a Frobenius fibre is $k[x,y]/(x^2,y^2)$, a quotient of $k[x,y]/(x,y)^3$. Lemma~\ref{lem:theta-jets} supplies the required evaluation in the smooth-theta case. The preceding product calculation supplies it in the other case. Nakayama's lemma then proves global generation.
\end{proof}

\section{The Cartier tensor and a relative determinant}
This appendix proves the cubic principal seed that begins the norm construction.


Put $L=\OO_A(3\Theta)$ and $W=H^0(A,L)$, of dimension nine. Fix invariant volume forms on $A$ and $A^{(2)}$. Relative Frobenius duality and $F^*L^{(2)}=L^2$ give a pairing
\begin{equation}\label{eq:beta}
 \beta:W\otimes W\longrightarrow W^{(2)},\qquad
 (s,t)\longmapsto\Tr_F(st).
\end{equation}
The trace here is the trace on top forms, transported by the chosen trivializations, not the field trace of a purely inseparable extension. Locally it is the coefficient-of-$xy$ map on a two-element $2$-basis. In particular the sheaf pairing
\begin{equation}\label{eq:perfect}
 F_*L\otimes F_*L\longrightarrow L^{(2)}
\end{equation}
is perfect: in the basis $1,x,y,xy$, each monomial pairs nontrivially with its complementary monomial. This is the explicit finite-Frobenius duality of \citet[Lemma~1.3.6, Proposition~1.3.7 and Theorem~1.3.8, pp.~26--28]{BrionKumar}.

\begin{lemma}[The typed symmetric tensor]\label{lem:tensor}
There is a basis $(e_a)_{a\in\FF_3^2}$ of $W$ identifying $W^{(2)}$ with $W^*$ under the finite level-three Heisenberg group such that the trilinear tensor associated to \eqref{eq:beta} is fully symmetric. It satisfies
\[
 T(e_a,e_a,e_b)=c\,\delta_{ab},
\]
where $c=0$ exactly when $A$ is nonordinary. In the nonordinary case $T$ is alternating and is specified by four scalars $\alpha_i$, one for each direction $i\in\PP^1(\FF_3)$: its value on three distinct basis vectors is $\alpha_i$ when their indices lie on a line of direction $i$, and is zero otherwise.
\end{lemma}
\begin{proof}
The theta group of $L$ has quotient $A[3]$ and admits a theta structure. Its weight-one representation on $W$ is the Schr\"odinger representation, by \citet[Lemma~8.24, p.~121; Theorem~8.28 and description~8.29, pp.~122--123; Corollary~9.17, pp.~133--134]{EGM}. Fix $\zeta\in\mu_3(k)\setminus\{1\}$. The finite subgroup generated by
\[
 e_a\longmapsto e_{a+x},\qquad
 e_a\longmapsto\zeta^{y\cdot a}e_a,\qquad
 e_a\longmapsto\zeta e_a
\]
has order $3^5$. If $\rho$ denotes its matrices, the Frobenius-twisted matrices obey $\rho(h)^{[2]}=\rho(h)^{-t}$. This holds directly for each displayed generator, since $\zeta^2=\zeta^{-1}$ and the translation matrices are orthogonal. Thus the indicated basis pairing identifies $W^{(2)}$ with the contragredient representation. This statement is restricted to the finite central $\mu_3$; it is not an assertion about the whole scalar $\mathbf G_m$.

Functoriality of the Frobenius trace and projection formula make \eqref{eq:beta} equivariant for this twisted action. Accordingly $T$ is invariant under the finite group and symmetric in its first two arguments. Modulations force a nonzero coefficient to have indices $a+b+c=0$. Translation by $b-a$ sends such a triple to $(b,c,a)$. Cyclic symmetry, together with the first transposition, gives full symmetry. If two indices coincide, the sum-zero condition forces all three to coincide. Projection formula gives $\beta(s,s)=c s^{(2)}$, where $c$ is the coefficient of Cartier on the invariant top form. Hence the diagonal coefficients are as asserted. This scalar is nonzero exactly when $A$ is Frobenius split: by finite Frobenius duality and $\omega_A\simeq\OO_A$, the space of maps $F_*\OO_A\to\OO_{A^{(2)}}$ is one-dimensional, and evaluation at $1$ is multiplication by $c$. A nonzero evaluation can be rescaled to split the unit map. Frobenius splitting is equivalent to ordinarity for an abelian variety \citep[Definition--Proposition~2.1, p.~1071]{SannaiTanaka}. When it is zero the tensor is alternating. Every triple of distinct indices with sum zero is an affine line, and translations identify coefficients on parallel lines, giving the four parameters.
\end{proof}

The following cubic is the reduction modulo two of the known Coble-cubic formula of \citet[equation~(3.9), author version p.~14]{GrusonSam}. We give its polynomial proof and characteristic-two exceptional-tensor analysis in full. For $I=\FF_3^2$, let
\[
 P_*(s)=\sum_{a\in I}s_a^3,\qquad
 P_i(s)=\sum_{\ell\parallel i}\prod_{a\in\ell}s_a.
\]
The second sum has three terms, one for each affine line in direction $i$.

\begin{lemma}[Pfaffian identity and its exceptional tensors]\label{lem:pfaffian}
Let $T$ be the alternating tensor with the four parameters of Lemma~\ref{lem:tensor}, and set $M(s)_{ab}=T(s,e_b,e_a)$. Then
\begin{equation}\label{eq:q}
 q_\alpha(s)=\left(\prod_i\alpha_i\right)P_*(s)
 +\sum_i\alpha_i\left(\sum_{j\ne i}\alpha_j^3\right)P_i(s)
\end{equation}
satisfies
\[
 M(s)s=0,\qquad
 \operatorname{Pf}(M(s)_{\widehat a,\widehat a})=q_\alpha(s)s_a,
 \qquad\operatorname{adj}M(s)=q_\alpha(s)^2ss^t.
\]
For $\alpha\ne0$, $q_\alpha$ is identically zero precisely at forty projective parameter points: the four points with one nonzero coordinate, and the thirty-six points with three nonzero coordinates whose cubes are equal. At every such point, $W$ is an orthogonal direct sum of three three-dimensional spaces on each of which $T$ is a nonzero volume form.
\end{lemma}
\begin{proof}
Perform the calculation over $\FF_2[\alpha_i,s_a]$. Alternation gives $Ms=0$. The vector $v$ of principal eight-dimensional Pfaffians is also in the kernel, by the Pfaffian expansion identity. At $s=e_0$, the minor omitting zero consists of four alternating two-by-two blocks, one for each direction; its Pfaffian is $\prod_i\alpha_i$. Thus the universal matrix has rank eight, so $v=qs$ in the fraction field. The polynomial identities $v_as_b=v_bs_a$ and the pairwise coprimality of the variables $s_a$ imply that $s_a$ divides $v_a$. Therefore $q$ is a polynomial, of degree three in $s$ and four in $\alpha$.

The finite Heisenberg generators preserve $T$ and have determinant one. The transformation rule for $\bigwedge^8W^*\simeq W\otimes\det W^*$ shows that $v$ transforms as $s$, and hence $q$ is invariant. A modulation-invariant cubic monomial has its three indices summing to zero. It is either $s_a^3$ or the product on an affine line. Translation invariance therefore puts $q$ in the span of $P_*,P_i$. The calculation at $e_0$ gives its $P_*$ coefficient.

For a line coefficient, set all section coordinates outside the vertical line $x=0$ equal to zero, and denote its three remaining coordinates by $x_0,x_1,x_2$. The minor omitting $(0,0)$ has a two-by-two block of entry $\alpha_\infty x_0$ and an off-diagonal six-by-six block whose Pfaffian is the determinant of
\[
 B=\begin{pmatrix}
 \alpha_0x_0&\alpha_1x_2&\alpha_2x_1\\
 \alpha_2x_2&\alpha_0x_1&\alpha_1x_0\\
 \alpha_1x_1&\alpha_2x_0&\alpha_0x_2
 \end{pmatrix}.
\]
Expansion of its six terms gives
\[
 \det B=\alpha_0\alpha_1\alpha_2(x_0^3+x_1^3+x_2^3)
 +(\alpha_0^3+\alpha_1^3+\alpha_2^3)x_0x_1x_2.
\]
The principal Pfaffian is $\alpha_\infty x_0\det B$; dividing that Pfaffian by $x_0$ gives $q|_{x=0}=\alpha_\infty\det B$, and hence the coefficient of $P_\infty$ in \eqref{eq:q}. Linear permutations of the index plane give the other three coefficients. The identity $\operatorname{adj}M=vv^t$ for odd alternating matrices then gives the adjugate formula.

The five orbit sums have disjoint supports. Consequently $q=0$ is equivalent to
\[
 \prod_i\alpha_i=0,\qquad
 \alpha_i\sum_{j\ne i}\alpha_j^3=0\quad\text{for each }i.
\]
If $r$ coordinates are nonzero, these equations make all their cubes equal to a common $u\ne0$ and require $ru=u$ in characteristic two. Thus $r$ is odd, and the product condition leaves $r=1$ or $3$. The projective counts are $4$ and $4\cdot3^2$.

When $r=1$, the three parallel lines partition the basis and give the required decomposition. When $r=3$, rescale so that all nonzero coefficients lie in $\mu_3$. A basis change $e_a\mapsto\zeta^{Q(a)}e_a$ with $Q$ a binary quadratic form over $\FF_3$ multiplies the coefficient in direction $v$ by $\zeta^{-Q(v)}$, because $\sum_{t\in\FF_3}Q(a+tv)=-Q(v)$. Evaluation of binary quadratics at three distinct directions is an isomorphism, so all three coefficients can be made one. Move the missing direction to the vertical direction by a linear change of the index plane.

For the final basis change, start with the sum of the three volume forms on horizontal lines and set
\[
 f_{x,u}=\sum_{y\in\FF_3}\zeta^{uy}e_{x,y}.
\]
This Fourier matrix is invertible; its inverse uses $\zeta^{-uy}$, since $3=1$ in $k$. Its transformed tensor is one exactly when the three first indices are distinct and the second indices sum to zero, and zero otherwise. These are precisely the lines of the three nonvertical directions. Thus the normalized three-direction tensor is equivalent to the one-direction tensor. Transporting its three orthogonal summands proves the last assertion.
\end{proof}

\begin{proposition}[The cubic principal seed]\label{prop:seed}
For every principally polarized abelian surface, a generic smooth curve in $|3\Theta|$ has genus ten and $2$-rank $8+f_A$. Its relative Frobenius is bijective on its eight-dimensional quotient.
\end{proposition}
\begin{proof}
The pairing \eqref{eq:perfect} and Corollary~\ref{cor:fgenerated} imply that $\beta$ is nonzero. Suppose first that $A$ is nonordinary. If $q_\alpha=0$, Lemma~\ref{lem:pfaffian} gives a rank-three projector $D$ on $W$ such that
\[
 T(Ds,t,u)=T(s,Dt,u)
\]
and whose summands are mutually orthogonal. It preserves the kernel of every fibre evaluation $W\to(F_*L)_x$. Indeed, if $a+b$ vanishes in the fibre with $a$ in one summand and $b$ in its orthogonal complement, the value of $a$ pairs to zero with the complement by orthogonality and with its own summand because $a+b$ has zero value. Global generation and perfectness force $a$ to have zero value. Thus the kernel bundle of $W\otimes\OO\twoheadrightarrow F_*L$ is preserved by $D$: the induced map from it to $F_*L$ vanishes at every closed point and hence identically. The projector descends to a nontrivial endomorphism of $F_*L$.

Oda's theorem contradicts this. The restriction of $\lambda_L=3\lambda_\Theta$ to the scheme-theoretic kernel of $F$ is a closed immersion, because $\lambda_\Theta$ is an isomorphism and multiplication by three is an automorphism on $A[2]$. Hence $\End(F_*L)=k$ by \citet[Theorem~1.2(i), p.~45; proof pp.~43--45]{Oda}. A scalar cannot act on $W$ as a rank-three projector. Therefore $q_\alpha\ne0$.

It remains essential to compute the map on the relative cohomology. By Lemma~\ref{lem:cohomology} and Serre duality, the fibre is the hyperplane $\ann(s)\subset W^*$, and its Frobenius-domain hyperplane is $\ann(s^{[2]})$. The contraction matrix $M(s)$ is precisely the linearization of Frobenius: duality converts the multiplication by the defining section in the restriction sequence to \eqref{eq:beta}. Since $Ms=0$, it maps the indicated domain into the indicated target. On $s_0=1$, the domain coordinate basis is $e_j^*+s_j^2e_0^*$ and the target basis is $e_i^*+s_ie_0^*$, for $1\le i,j\le8$. The relative matrix is
\[
 R(s)_{ij}=M(s)_{ij}+s_j^2M(s)_{i0}.
\]
Expansion by the rank-one column additions, using the adjugate in Lemma~\ref{lem:pfaffian}, gives
\begin{equation}\label{eq:relative-det}
 \det R(s)=q_\alpha(s)^2\sum_{a\in I}s_a^3.
\end{equation}
For completeness, the unmodified minor contributes $q^2s_0^2=q^2$. Replacing its $j$th column by $s_j^2(M_{i0})_{i\ne0}$ contributes $s_j^2q^2s_0s_j=q^2s_j^3$; terms with two replaced columns vanish. This proves \eqref{eq:relative-det}. Both factors are nonzero polynomials. Thus the relative map is generically invertible, which proves stable rank eight rather than only an initial rank assertion.

If $A$ is ordinary, Lemma~\ref{lem:tensor} gives a symmetric contraction matrix with diagonal $c s_a$, where $c\ne0$. In characteristic two, terms of a symmetric determinant with a permutation cycle of length at least three cancel with their inverse-cycle partners. The remaining terms are products of diagonal entries and squares of off-diagonal entries. The monomial $\prod_as_a$ therefore has coefficient $c^9\ne0$, since every nondiagonal term has some positive even exponent. The full matrix is generically invertible. It preserves the relative hyperplanes: full symmetry gives $M(s)s=c s^{[2]}$. Therefore its restriction between these equal-dimensional hyperplanes is also invertible. This proves the same relative assertion.

Finally $|3\Theta|$ is very ample, so its smooth-member open meets the nonvanishing determinant open. Equation~\eqref{eq:dimensions} gives genus ten and $2$-rank $8+f_A$.
\end{proof}


\section{Nodal norms and their linear systems}

This appendix proves nodality of general norm images and places their translates in the required complete linear systems.

\begin{lemma}[Nodal norm images]\label{lem:norm-nodes}
Let $\phi:B\to A$ be a finite \'etale isogeny, where $(B,\Theta)$ is principally polarized. For every $n\ge3$, a general smooth $C\in|n\Theta|$ maps birationally to an integral nodal curve, with exactly $n^2(\deg\phi-1)$ nodes. This assertion allows even $\deg\phi$.
\end{lemma}
\begin{proof}
Let $K=\ker\phi(k)$, a finite constant group. We exclude tangencies and triple identifications for the finitely many nonzero translations in $K$.

Suppose first that $\Theta$ is smooth. For $n=3$, Lemma~\ref{lem:theta-jets} shows that the six first-jet coordinates at $P,P+g$ are independent outside a finite set of $P$. For $n\ge4$ they are independent everywhere. At such a pair, requiring the two values to vanish gives two independent equations, and requiring proportional gradients gives a further nonzero determinant equation. Thus the tangency incidence has codimension three in the section space while $P$ varies in dimension two. At the finite exceptional set for $n=3$, the two independent value equations suffice to avoid the pair altogether. For distinct $0,g,h\in K$, three-value separation imposes three conditions at $P,P+g,P+h$ while $P$ varies in dimension two. Hence triple identifications are avoided on a nonempty open.

We spell out the product case, where cubic elliptic factors have a smaller jet space. A degree-$n$ line bundle on an elliptic curve separates two values and separates a first jet at a single point for $n\ge3$. For $n=3$, the evaluation on the two first jets at distinct $x,y$ has rank three; after their two values are prescribed to be zero, the two derivatives form a line. Both derivative coefficients on this line are nonzero unless the unique residual zero equals $x$ or $y$. For $y=x+g$, these exceptions are the finite equations $3x+g=0$ or $3x+2g=0$, after the chosen origin is incorporated in the line bundle class.

When $g=(g_1,g_2)$ has both coordinates nonzero, at points avoiding those finite exceptions in both factors the tensor-product section space supplies the full six surface first jets. To see this, split each three-dimensional factor space into its two value coordinates and its one-dimensional common derivative direction. The tensor of the two value spaces supplies both surface values; derivative-direction times value-space supplies independent derivatives at the two points in one direction, and value-space times derivative-direction does the same in the other direction. It also suffices that just one factor have both derivative coefficients nonzero. To verify this, use value bases $u_1,u_2$ and $v_1,v_2$ and residual derivative directions $u_3,v_3$. Suppose the second factor has both coefficients nonzero. Modulo the two surface values, the two first-direction derivatives involve respectively the disjoint pairs of coordinates $(u_3v_1,u_2v_1)$ and $(u_3v_2,u_1v_2)$, with each pair having a nonzero coefficient because each individual elliptic first jet is onto. The second-direction derivatives can then be adjusted independently by $u_1v_3$ and $u_2v_3$. All four derivatives are therefore independent. The locus where both factors fail is finite and is avoided by the two independent value conditions.

If precisely one $g_i$ is nonzero, the cubic tensor evaluation has rank five. After the two values vanish, the two derivatives in the varying coordinate lie on a nonzero line, and the two transverse derivatives are independent. Their tangency determinant is a nonzero polynomial in these three parameters. This again gives three conditions. The possible further derivative exceptions occur on a one-dimensional locus of $P$; there the same nonzero derivative line and independent transverse derivatives still give the third condition. Thus its incidence is proper as well.

For triples in a cubic product, three Segre value tensors are dependent only if one coordinate is constant and the three value vectors in the other plane cubic are collinear. Indeed, if two vectors in each factor are independent, their two simple tensors span no third simple tensor except on the two ruling lines; this follows immediately in the resulting two-by-two tensor matrix. A plane cubic triple $x,x+g,x+h$ is collinear only when $3x+g+h=0$. Hence the rank-two exceptional locus has dimension at most one and the two value conditions exclude it generically. Elsewhere the triple values have rank three and the point parameter has dimension two.

For $n\ge4$ on a product, three distinct elliptic values are independent. For $n\ge5$, two elliptic first jets are independent as well. For $n=4$, their rank can drop from four to three only if the line bundle is $\OO(2x+2y)$; for $y=x+g$ this is a finite condition $4x+2g=0$. Tensor products therefore give six surface first jets away from a finite set when both $g_i$ are nonzero. When one $g_i$ vanishes, the exceptional locus has dimension one and rank five, with the same nonzero tangency determinant described above. Three distinct surface values are independent: separate the first-coordinate classes and then the at most three second-coordinate values within each class. These assertions follow in each elliptic factor from Riemann--Roch applied to the evaluation kernel.

All the undesirable incidences are thus proper. A general section also has no nontrivial kernel stabilizer. If translation by $g$ does not preserve the line bundle it preserves no member; if it does, its projective action in the very ample embedding is nontrivial, and its invariant members form a proper closed subset. Therefore $C\to\phi(C)$ has generic degree one. Since $\phi$ is \'etale, the only singularities of the image are the transverse pairs just obtained, so they are nodes. For each nonzero $g$, the curves $C$ and $C+g$ meet in $2n^2$ points. These are ordered branch pairs; every node contributes two. Consequently the number of nodes is $n^2(|K|-1)$. This is a count of finite sets and does not divide by two in the coefficient field.
\end{proof}

\begin{lemma}\label{lem:exact-system}
The norm of a very ample line bundle under a finite \'etale morphism of projective varieties over $k$ is very ample. On an abelian variety, any effective divisor in the numerical class of an ample $N$ can be translated into the exact system $|N|$.
\end{lemma}
\begin{proof}
To separate two target points by a norm section, choose a source section vanishing at one selected point over the first and nonzero at all the other finitely many points over both targets. Very ampleness makes each unwanted vanishing a proper hyperplane in the space of sections vanishing at the selected point. To separate a tangent vector, choose its lifted derivative nonzero at the selected point and keep all other fibre values nonzero. The norm is locally a product; its other factors are units, so its derivative is the selected nonzero derivative times a unit. Basepoint freeness is proved by a section avoiding the entire finite fibre. Thus norm sections themselves separate points and tangent vectors.

For the second assertion, the two line bundles differ by an element of $\Pic^0$. Surjectivity of $\lambda_N$ realizes that difference by a translation of $N$, and translating the divisor gives the required member. Translations preserve smoothness, nodes and cohomological ranks.
\end{proof}

\section*{Sources and method}
This section records the primary inputs, computational scope and verification method.
The integral and cover-theoretic results used here are cited at primary-source locations. The principal cubic and nodal-norm arguments in the appendices also occur in the companion treatment of principal and odd-degree polarizations; they are included here for self-contained verification. The accompanying program \emph{Local algebra certificate} checks the universal factorization, transition and residual-complex identities by independent sparse-polynomial arithmetic. Its scope is local algebra, not global determinant nonvanishing. Source code and records are available at \url{\workbenchurl}; permanent records are at \url{\recordurl}.

The mathematics and text were produced by GPT-6 Astra (ChatGPT Pro mode). Separate same-model reconstructions from statements alone establish the rank theorems, maximal-root reduction and conditional determinant theorem; their full comparisons found no mathematical defect. The rank reconstruction uses precisely stated principal-polarization inputs and a toral Cartier upper bound, while the Prym-specific auxiliary proof has a full targeted check. The ample-difference theorem has a separate same-model reconstruction from its statement and the exact index-zero and positive-index inputs; its complete comparison also passed. Its numerical ampleness proof uses the cited Riemann--Roch and finite-kernel criteria. The finite-field existence and eventual-rank theorem has a separate same-model statement-led reconstruction relative to the exact integral-root and ample-difference inputs, followed by reciprocal full comparisons. Its explicit threshold was separately reconstructed from its statement before comparison. These reconstructions use the earlier polarization, integral-root and smoothing results as stated inputs. The eventual theorem for arbitrary mixed pairs has a separate same-model reconstruction from its statement, using the finite-field and exact-rank theorems as inputs, followed by reciprocal full comparisons. The no-index-zero-class theorem and the elementary $\alpha_2$ descent lemma also have separate same-model reconstructions and full comparisons, using the cited geometric-monodromy and abelian-variety inputs. Their detailed source and exposure scopes are recorded separately. The determinant reconstruction includes the torsor-cup comparison, the Frobenius-typed variation, the restricted direction space and the full residual boundary kernel. Its local frame calculations use faithfully flat descent, including parameter derivatives. The auxiliary smooth-jet and detailed norm-pencil assertions have targeted checks, without a separate-reconstruction claim. The automorphism contribution (P11)--(P12) has a separate statement-led derivation and full targeted comparison. Uniform determinant nonvanishing remains the open question stated above. The complete text received expert and newcomer readings by the same model. Reasoning-effort metadata is not exposed. No human has checked the mathematics. The new text and code are dedicated to the public domain under CC0 1.0.

\begin{thebibliography}{99}
\bibitem[Achter and Pries(2008)]{AchterPries}
Achter, J. D., \& Pries, R. (2008). Monodromy of the $p$-rank strata of the moduli space of curves. \emph{International Mathematics Research Notices, 2008}(15), rnn053. \url{https://doi.org/10.1093/imrn/rnn053}. Author version: \url{https://www.math.colostate.edu/~achter/math/prm0508.pdf}.
\bibitem[Brion and Kumar(2005)]{BrionKumar}
Brion, M., \& Kumar, S. (2005). \emph{Frobenius splitting methods in geometry and representation theory}. Progress in Mathematics, 231. Birkh\"auser. \url{https://www-fourier.univ-grenoble-alpes.fr/~mbrion/fbook.pdf}.
\bibitem[Chai(2005)]{Chai}
Chai, C.-L. (2005). Monodromy of Hecke-invariant subvarieties. \emph{Pure and Applied Mathematics Quarterly, 1}(2), 291--303. \url{https://doi.org/10.4310/PAMQ.2005.v1.n2.a4}.
\bibitem[Crew(1984)]{Crew}
Crew, R. M. (1984). Etale $p$-covers in characteristic $p$. \emph{Compositio Mathematica, 52}(1), 31--45. \url{https://www.numdam.org/item/CM_1984__52_1_31_0/}.
\bibitem[Debarre(1994)]{Debarre}
Debarre, O. (1994). Degrees of curves in abelian varieties. \emph{Bulletin de la Soci\'et\'e Math\'ematique de France, 122}(3), 343--361. \url{https://www.math.ens.psl.eu/~debarre/25.pdf}.
\bibitem[Edixhoven et~al.(n.d.)]{EGM}
Edixhoven, B., van der Geer, G., \& Moonen, B. (n.d.). \emph{Abelian varieties: Preliminary version of the first chapters}. Author manuscript. \url{https://van-der-geer.nl/~gerard/AV.pdf}.
\bibitem[Genestier and Ng{\^o}(2009)]{GenestierNgo}
Genestier, A., \& Ng{\^o}, B. C. (2009). Lectures on Shimura varieties. \emph{Panoramas et Synth\`eses, 29}, 187--236. Author version. \url{https://www.math.uchicago.edu/~ngo/Shimura.pdf}.
\bibitem[Grothendieck(n.d.)]{Grothendieck}
Grothendieck, A. (n.d.). \emph{Notes on crystalline cohomology}, cote~4, \S8.5.8, archive sheets~7--8. Institut Montpelli\'erain Alexander Grothendieck. \url{https://grothendieck.umontpellier.fr/4.pdf}.
\bibitem[Gruson and Sam(2015)]{GrusonSam}
Gruson, L., \& Sam, S. V. (2015). Alternating trilinear forms on a nine-dimensional space and degenerations of $(3,3)$-polarized abelian surfaces. \emph{Proceedings of the London Mathematical Society, 110}(3), 755--785. \url{https://doi.org/10.1112/plms/pdu050}. Author version: \url{https://arxiv.org/abs/1301.5276}.
\bibitem[Oda(1969)]{Oda1969}
Oda, T. (1969). The first de Rham cohomology group and Dieudonn\'e modules. \emph{Annales scientifiques de l'\'Ecole normale sup\'erieure, 2}(1), 63--135. \href{https://doi.org/10.24033/asens.1175}{DOI: 10.24033/asens.1175}.
\bibitem[Oda(1971)]{Oda}
Oda, T. (1971). Vector bundles on an elliptic curve. \emph{Nagoya Mathematical Journal, 43}, 41--72. \url{https://doi.org/10.1017/S0027763000014367}.
\bibitem[Pries(2008)]{PriesTorsion}
Pries, R. (2008). A short guide to $p$-torsion of abelian varieties in characteristic $p$. In \emph{Computational arithmetic geometry} (Contemporary Mathematics, Vol.~463, pp.~121--129). American Mathematical Society. \url{https://arxiv.org/abs/math/0609658}.
\bibitem[Pries and Zhu(2012)]{PriesZhu}
Pries, R., \& Zhu, H. J. (2012). The $p$-rank stratification of Artin--Schreier curves. \emph{Annales de l'Institut Fourier, 62}(2), 707--726. \url{https://doi.org/10.5802/aif.2692}.

\bibitem[Sannai and Tanaka(2016)]{SannaiTanaka}
Sannai, A., \& Tanaka, H. (2016). A characterization of ordinary abelian varieties by the Frobenius push-forward of the structure sheaf. \emph{Mathematische Annalen, 366}, 1067--1087. \url{https://doi.org/10.1007/s00208-015-1352-3}.
\bibitem[The Stacks Project Authors(2026)]{Stacks}
The Stacks Project Authors. (2026). \emph{The Stacks project}. Tags~01ZM, 0B8W, 0C0C, 081F, 09MT, 0B91, 0B9T, 02O1, 0D4E and 062Y. \url{https://stacks.math.columbia.edu/}.
\bibitem[Zou(2026)]{Zou}
Zou, H. (2026, September 29). \emph{Quasi-$F$-split primitive symplectic varieties in positive characteristic}. Author preprint. \url{https://www.math.uni-bielefeld.de/~hzou/papers/qfpsv/}.
\end{thebibliography}
\end{document}
